%% 第 5 章 Hecke 特征 D-模（原书 p.192--p.229）。
%% 2026-11-21 校订时把 units/ch05_u01--u04 内联展开为本文件；此后本章不再由 assemble.py 覆盖。

% ===== units/ch05_u01.tex =====
% ch05_u01 —— 原书 p.192–p.200（第 5 章开头：5.1 全部；5.2 全部；5.3 自 5.3.1 起；p.200 末句中断于「So by」）
\bdchap{ch:hecke}{Hecke 特征 D-模}
\SourcePage{192}

\bdsec{5.1. D-模的构造}

\bdp{5.1.1} 本小节在 $\operatorname{Bun}_G$ 上构造一族 $\mathcal{D}$-模，以 $\mathcal{O}\mathfrak{p}_{{}^{L}G}(X)$ 为参数；后者即 $X$ 上 ${}^{L}G$-oper 的叠。

记 $Z$ 为 ${}^{L}G$ 的中心。按 3.4.3 的公式 (57)，须对 $\mathcal{L}\in Z\operatorname{tors}_\theta(X)$ 配置 $\operatorname{Bun}_G$ 上一族 $\mathcal{D}$-模，以 $\mathcal{O}\mathfrak{p}_{{}^{L}\mathfrak{g}}(X)$ 为参数。在 4.4.3 中我们定义了 $\lambda_{\mathcal{L}}\in\mu_\infty\operatorname{tors}_\theta(\operatorname{Bun}_G)$。$\lambda_{\mathcal{L}}$ 是 $\operatorname{Bun}_G$ 上的线丛，带有同构 $\lambda_{\mathcal{L}}^{\otimes 2n}\xrightarrow{\sim}(\omega^{\sharp}_{\operatorname{Bun}_G})^{\otimes n}$，其中 $n\ne 0$（见 4.0.1）。故 $\lambda_{\mathcal{L}}$ 是 $\mathcal{D}'$-模。因此 $M_{\mathcal{L}}:=\lambda_{\mathcal{L}}^{-1}\otimes_{\mathcal{O}_{\operatorname{Bun}_G}}\mathcal{D}'$ 是 $\operatorname{Bun}_G$ 上的左 $\mathcal{D}$-模。按 3.3.2 与 2.7.4，有典范的代数态射 $h_X\varphi_X:A_{{}^{L}\mathfrak{g}}(X)\to\Gamma(\operatorname{Bun}_G,\mathcal{D}')$。于是 $\Gamma(\operatorname{Bun}_G,\mathcal{D}')$ 在 $\mathcal{D}'$ 上的右作用，给出 $M_{\mathcal{L}}$ 上的 $A_{{}^{L}\mathfrak{g}}(X)$-模结构。因此可以把 $M_{\mathcal{L}}$ 看作 $\operatorname{Bun}_G$ 上一族左 $\mathcal{D}$-模，以 $\operatorname{Spec}A_{{}^{L}\mathfrak{g}}(X)=\mathcal{O}\mathfrak{p}_{{}^{L}\mathfrak{g}}(X)$ 为参数。
% [Tutor Note: 原书 p.192 把复合态射印作 $h_X\varphi_X$（$h_X$ 与 $\varphi_X$ 之间无复合记号），底本如此，保留原样。]

这样，我们在 $\operatorname{Bun}_G$ 上构造了一族左 $\mathcal{D}$-模，以 $\mathcal{O}\mathfrak{p}_{{}^{L}G}(X)$ 为参数。对 ${}^{L}G$-oper $\mathfrak{F}$，相应的 $\mathcal{D}$-模 $M_{\mathfrak{F}}$ 是 $M_{\mathcal{L}}/\mathfrak{m}_{\mathfrak{F}}M_{\mathcal{L}}=\lambda_{\mathcal{L}}^{-1}\otimes\mathcal{D}'/\mathcal{D}'\mathfrak{m}_{\mathfrak{F}}$，其中 $\mathcal{L}$ 是 $\mathfrak{F}$ 在 $Z\operatorname{tors}_\theta(X)$ 中的像，而 $\mathfrak{m}_{\mathfrak{F}}\subset A_{{}^{L}\mathfrak{g}}(X)$ 是对应于 $\mathfrak{F}$ 的 ${}^{L}\mathfrak{g}$-oper 的极大理想。

\bdp{5.1.2} \emph{命题.}
\begin{enumerate}[label=(\roman*),leftmargin=2.2em]
\item 对每个 $\mathcal{L}\in Z\operatorname{tors}_\theta(X)$，$M_{\mathcal{L}}$ 在 $A_{{}^{L}\mathfrak{g}}(X)$ 上平坦。
\item 对每个 ${}^{L}G$-oper $\mathfrak{F}$，$\mathcal{D}$-模 $M_{\mathfrak{F}}$ 完整；其奇异支集作为闭链，与 Hitchin 纤维化的零纤维重合。
\end{enumerate}

\emph{证明.} 按 2.2.4(iii)，$\operatorname{gr}\mathcal{D}'$ 在 $\operatorname{gr}A_{{}^{L}\mathfrak{g}}(X)$ 上平坦\footnote{这意指：若 $f:S\to\operatorname{Bun}_G$ 光滑且 $S$ 仿射，则 $\Gamma(S,f^*\operatorname{gr}\mathcal{D}')$ 是自由 $\operatorname{gr}A_{{}^{L}\mathfrak{g}}(X)$-模（在 $A_0=\mathbb{C}$ 的 $\mathbb{Z}_+$-分次环 $A$ 上，平坦的 $\mathbb{Z}_+$-分次模是自由的）。}。故 $\mathcal{D}'$ 在 $A_{{}^{L}\mathfrak{g}}(X)$ 上平坦。这蕴含 i)，以及对任意理想 $I\subset A_{{}^{L}\mathfrak{g}}(X)$ 的等式 $\operatorname{gr}(\mathcal{D}'/\mathcal{D}'I)=\operatorname{gr}\mathcal{D}'/(\operatorname{gr}\mathcal{D}'\cdot\operatorname{gr}I)$。若 $I$ 极大，即得 ii)。\hfill$\Box$

\SourcePage{193}

\bdsec{5.2. 主要定理 I：导引}

\bdp{5.2.1} 整体主定理 5.2.6 断言：$\mathcal{D}$-模 $M_{\mathfrak{F}}$ 是\emph{Hecke 函子}的特征模。为定义 Hecke 函子，我们引入\emph{大 Hecke 叠} $\mathcal{H}ecke$。$S$-点构成的群胚 $\mathcal{H}ecke(S)$ 由四元组 $(\mathcal{F}_1,\mathcal{F}_2,x,\alpha)$ 组成，其中 $\mathcal{F}_1,\mathcal{F}_2$ 是 $X\times S$ 上的 $G$-torsor，$x\in X(S)$，而 $\alpha:\mathcal{F}_1|_U\xrightarrow{\sim}\mathcal{F}_2|_U$ 是 $x$ 的图像之补集 $U$ 上的同构。有显然的投影 $p_{1,2,X}=(p_1,p_2,p_X):\mathcal{H}ecke\to\operatorname{Bun}_G\times\operatorname{Bun}_G\times X$。

叠 $\mathcal{H}ecke$ 是归纳代数的，投影 $p_i$、$p_{i,X}$ 是归纳固有的。确切地说，存在闭代数子叠的递增族 $\mathcal{H}ecke_1\subset\mathcal{H}ecke_2\subset\cdots\subset\mathcal{H}ecke$，使得 $\mathcal{H}ecke=\bigcup\mathcal{H}ecke_a$，且 $p_i:\mathcal{H}ecke_a\to\operatorname{Bun}_G$ 与 $p_{i,X}:\mathcal{H}ecke_a\to\operatorname{Bun}_G\times X$ 都是固有态射。

\bdp{5.2.2} \emph{注.} (i) 诸 $\alpha$ 的复合使 $\mathcal{H}ecke$ 成为 $\operatorname{Bun}_G$ 上的 $X$-群胚族。

(ii) $\mathcal{H}ecke$ 是 $\operatorname{Bun}_G\times X$ 上扭仿射 Grassmann 流形的族。确切地说，对 $(\mathcal{F}_2,x)\in\operatorname{Bun}_G\times X$，纤维 $\mathcal{H}ecke_{(\mathcal{F}_2,x)}:=p_{2,X}^{-1}(\mathcal{F}_2,x)$ 典范同构于仿射 Grassmann 流形 $\mathcal{GR}_x:=G(K_x)/G(O_x)$ 经 $G(O_x)$-torsor $\mathcal{F}_2(O_x)$ 所作的扭（就左 $G(O_x)$-作用而言）。当 $\mathcal{F}_2$ 是平凡丛时，我们在 4.5.2 中描述了这个同构。一般情形的构造类似：固定 $\gamma_2\in\mathcal{F}_2(O_x)$，把 $(\mathcal{F}_1,\mathcal{F}_2,x,\alpha)$ 对应到 $\gamma_2/\alpha(\gamma_1)$ 在 $G(K_x)/G(O_x)$ 中的像；这里 $\gamma_1$ 是 $\mathcal{F}_1(O_x)$ 的任一元素，$\gamma_2/\alpha(\gamma_1)$ 表示满足 $g\alpha(\gamma_1)=\gamma_2$ 的元素 $g\in G(K_x)$。按 2.3.4，态射 $\mathcal{H}ecke_{(\mathcal{F}_2,x)}\to G(K_x)/G(O_x)$ 是同构。

\bdp{5.2.3} 态射 $\nu:\mathbb{G}_m\to G$ 的共轭类集，可典范地等同于 ${}^{L}G$ 的支配权集 $P_+({}^{L}G)$。回忆 $\mathcal{GR}_x=G(K_x)/G(O_x)$ 中的 $G(O_x)$-轨道由 $\chi\in P_+({}^{L}G)$ 标号。按定义，$\operatorname{Orb}_\chi$ 是 $\nu(t_x)\in G(K_x)$ 在 $\mathcal{GR}_x$ 中的像的 $G(O_x)$-轨道，其中 $\nu:\mathbb{G}_m\to G$ 的类为 $\chi$，$t_x\in O_x$ 是素元。

\SourcePage{194}

按 5.2.2(ii)，$\mathcal{GR}_x$ 按 $\operatorname{Orb}_\chi$ 所作的分层，给出叠 $\mathcal{H}ecke$ 按子叠 $\mathcal{H}ecke_\chi$ 的分层，$\chi\in P_+({}^{L}G)$。$\mathcal{H}ecke_\chi$ 的 $\mathbb{C}$-点是四元组 $(\mathcal{F}_1,\mathcal{F}_2,x,\alpha)$，使得对某组 $\gamma_i\in\mathcal{F}_i(O_x)$ 与 $x$ 处某个形式参数 $t_x$，有 $\gamma_2=\nu(t_x)\alpha(\gamma_1)$，其中 $\nu:\mathbb{G}_m\to G$ 的类为 $\chi$。对合 $(\mathcal{F}_1,\mathcal{F}_2,x,\alpha)\mapsto(\mathcal{F}_2,\mathcal{F}_1,x,\alpha^{-1})$ 把 $\mathcal{H}ecke_\chi$ 与 $\mathcal{H}ecke_{\chi^\circ}$ 等同，其中 $\chi^\circ$ 是对偶权。于是 $p_{2,X}:\mathcal{H}ecke_\chi\to\operatorname{Bun}_G\times X$ 的纤维是 $\operatorname{Orb}_\chi$ 的扭形式，而 $p_{1,X}:\mathcal{H}ecke_\chi\to\operatorname{Bun}_G\times X$ 的纤维是 $\operatorname{Orb}_{\chi^\circ}$ 的扭形式。

对每个 $\chi$，叠 $\mathcal{H}ecke_\chi$ 在 $\operatorname{Bun}_G\times X$ 上光滑。通常其闭包 $\overline{\mathcal{H}ecke_\chi}$ 不光滑。
% [Tutor Note: 原书 p.194 此处既写出「closure」一词，又用上划线记号 $\overline{\mathcal{H}ecke_\chi}$，语义重复；底本如此，照录。]

\emph{注.} (i) 按 4.5.12，$\overline{\mathcal{H}ecke_\chi}$ 是诸分层 $\mathcal{H}ecke_{\chi'}$（$\chi'\le\chi$）之并。

(ii) 若 $G=GL(n)$，则我们的分层标号与「自然」标号一致。确切地说，设 $V_1,V_2$ 是对应于 $\mathfrak{F}_1,\mathfrak{F}_2$ 的向量丛。则 $\mathcal{H}ecke_\chi$ 由一切这样的组 $(V_1,V_2,x,\alpha)$ 组成：在 $x$ 的形式邻域上适当选取 $V_i$ 的基之后，$\alpha$ 的矩阵等于 $t_x^\chi$。

\bdp{5.2.4} 下面定义 Hecke 函子 $T_\chi^i:\mathcal{M}(\operatorname{Bun}_G)\to\mathcal{M}(\operatorname{Bun}_G\times X)$，其中 $\mathcal{M}$ 表示 $\mathcal{D}$-模范畴，$\chi\in P_+({}^{L}G)$，$i\in\mathbb{Z}$。

对 $\chi\in P_+({}^{L}G)$、$M\in\mathcal{M}(\operatorname{Bun}_G)$，记 $p_{1\chi}^\star M$ 为下述对象：$M$ 经光滑投影 $p_{1\chi}:\mathcal{H}ecke_\chi\to\operatorname{Bun}_G$ 的拉回，在 $\overline{\mathcal{H}ecke_\chi}$ 上的极小（$=$ Goresky--MacPherson）扩张；这里 $p_{1\chi}:=p_1|_{\mathcal{H}ecke_\chi}$。注意纤维化 $p_{1X}:\overline{\mathcal{H}ecke_\chi}\to\operatorname{Bun}_G\times X$ 局部平凡（见 5.2.2(ii)、5.2.3）。因此选定局部平凡化后，$p_{1\chi}^\star M$（局部地）等同于 $M$ 与仿射 Grassmann 流形上相应 $G(O)$-轨道\footnote{该轨道是 $\operatorname{Orb}_{\chi^\circ}$，其中 $\chi^\circ$ 是对偶权，见 5.2.3。}之闭包上的「相交上同调」$\mathcal{D}$-模的外张量积。

Hecke 函子 $T_\chi^i:\mathcal{M}(\operatorname{Bun}_G)\to\mathcal{M}(\operatorname{Bun}_G\times X)$ 定义为
\begin{equation}\tag{254}
T_\chi^i=H^i(p_{2,X})_*p_{1\chi}^\star
\end{equation}

\SourcePage{195}

其中 $H^i(p_{2,X})_*$ 是投影 $p_{2,X}:\overline{\mathcal{H}ecke_\chi}\to\operatorname{Bun}_G\times X$ 的上同调前推函子。

\emph{注.} 对局部有限型代数叠的可表拟紧态射 $f:\mathcal{X}\to\mathcal{Y}$，$H^if_*:\mathcal{M}(\mathcal{X})\to\mathcal{M}(\mathcal{Y})$ 的定义是清楚的。事实上，概形情形已有 $H^if_*$ 的定义，且已知 $H^if_*$ 与光滑基变换交换。

\bdp{5.2.5} 对 $\chi\in P_+({}^{L}G)$，记 $V^\chi$ 为最高权 $\chi$ 的不可约 ${}^{L}G$-模，并带一个标记的最高向量。若 $\mathfrak{F}$ 是 $X$ 上的 ${}^{L}G$-oper（更一般地，带联络的 ${}^{L}G$-丛），记 $V^\chi_{\mathfrak{F}}$ 为 $V^\chi$ 的 $\mathfrak{F}$-扭；它是 $X$ 上的光滑 $\mathcal{D}$-模。

\bdp{5.2.6} \emph{整体主定理.} 设 $\mathfrak{F}$ 是 $X$ 上的 ${}^{L}G$-oper，$M_{\mathfrak{F}}$ 是 5.1.1 中定义的 $\operatorname{Bun}_G$ 上的 $\mathcal{D}$-模。则当 $i\ne 0$ 时 $T_\chi^iM_{\mathfrak{F}}=0$，且有 $\operatorname{Bun}_G\times X$ 上 $\mathcal{D}$-模的典范同构
\begin{equation}\tag{255}
T_\chi^0M_{\mathfrak{F}}\xrightarrow{\sim}M_{\mathfrak{F}}\boxtimes V^\chi_{\mathfrak{F}}.
\end{equation}

同构 (255) 与 Hecke 对应的复合以及 $V^\chi$ 的张量积相容。确切陈述见 5.4.3。这一切意味着：$M_{\mathfrak{F}}$ 是以 $\mathfrak{F}$ 为特征值的\emph{Hecke 特征 $\mathcal{D}$-模}。

\bdp{5.2.7} Laumon 对 $X$ 上不可约局部系统 $E$ 定义了一个猜测中的「朗兰兹变换」$K_E$（见 [La87] 的 §5.3 与 4.3.3；当 $\operatorname{rank}E\le 2$ 时 $K_E$ 确实存在）。$K_E$ 是 $\operatorname{Bun}_{GL_n}$ 上的完整 $\mathcal{D}$-模，$n=\operatorname{rank}E$；至少当 $n=2$ 时，其奇异支集是 Hitchin 纤维化的零纤维（见 [La87] 的 §5.5）。此外 $K_E$ 有正则奇点，且它在 $\operatorname{Bun}_{GL_n}$ 每个连通分支上的限制不可约。若 $E$ 是 $SL_n$-局部系统，则 $K_E$ 位于 $\operatorname{Bun}_{PGL_n}$ 上。

考虑 5.1.2 与 5.2.6，自然可以猜测：当 $G=PGL_n$ 时，5.1.1 中的 $\mathcal{D}$-模 $M_{\mathfrak{F}}$ 等于 $K_{\mathfrak{F}}$（这方面的一些结果可见 [Fr]）。此外，还想弄清：\SourcePage{196}（对任意 $G$）$M_{\mathfrak{F}}$ 是否有正则奇点，它在 $\operatorname{Bun}_G$ 各连通分支上的限制是否不可约。

\bdp{5.2.8} 用 5.1.1 中的 $\mathcal{D}$-模 $M_{\mathcal{L}}$（$\mathcal{L}\in Z\operatorname{tors}_\theta(X)$）改写 5.2.6，既方便又重要。按 (57)，$\mathcal{L}\in Z\operatorname{tors}_\theta(X)$ 定义了 $X$ 上一族 ${}^{L}G$-oper $\mathfrak{F}_{\mathcal{L}}$，以 $\operatorname{Spec}A_{{}^{L}\mathfrak{g}}(X)$ 为参数。于是 $\mathfrak{F}_{\mathcal{L}}$ 是 $X\times\operatorname{Spec}A_{{}^{L}\mathfrak{g}}(X)$ 上的 ${}^{L}G$-torsor，带有沿 $X$ 的联络。对 $\chi\in P_+({}^{L}G)$，$V^\chi$ 的 $\mathfrak{F}_{\mathcal{L}}$-扭是 $X\times\operatorname{Spec}A_{{}^{L}\mathfrak{g}}(X)$ 上的向量丛，带有沿 $X$ 的联络。我们把它看作 $X$ 上的 $\mathcal{D}$-模 $V^\chi_{\mathcal{L}}$，并带有 $A_{{}^{L}\mathfrak{g}}(X)$ 的作用。

现在考虑 $\operatorname{Bun}_G$ 上的 $\mathcal{D}$-模 $M_{\mathcal{L}}$（见 5.1.1）；$A_{{}^{L}\mathfrak{g}}(X)$ 作用于它。容易看出（用 5.1.2(i)）：5.2.6 是下述定理的推论。

\bdp{5.2.9} \emph{定理.} 存在 $\operatorname{Bun}_G\times X$ 上 $\mathcal{D}$-模的典范同构
\begin{equation}\tag{256}
T_\chi^0M_{\mathcal{L}}\simeq M_{\mathcal{L}}\underset{A_{{}^{L}\mathfrak{g}}(X)}{\boxtimes}V^\chi_{\mathcal{L}},
\end{equation}
它与 $A_{{}^{L}\mathfrak{g}}(X)$ 的作用相容；并且当 $i\ne 0$ 时 $T_\chi^iM_{\mathcal{L}}=0$。

\bdp{5.2.10} 我们将从整体定理的局部版本推出它，现在就来说明这个局部版本。考虑仿射 Grassmann 流形 $\mathcal{GR}:=G(K)/G(O)$，其中 $O:=\mathbb{C}[[t]]$，$K=\mathbb{C}((t))$。这是归纳固有的归纳概形。于是有 $\mathcal{GR}$ 上 $\mathcal{D}$-模的「抽象」范畴 $\mathcal{M}(\mathcal{GR})$，它定义为 $\varinjlim\mathcal{M}(Y)$，其中 $Y$ 取遍 $\mathcal{GR}$ 的一切闭子概形。

我们无法把 $\mathcal{GR}$ 表示成光滑子概形递增序列之并。然而 $\mathcal{GR}$ 是\emph{形式光滑}的归纳概形。这允许把 $\mathcal{GR}$ 上的 $\mathcal{D}$-模当作「具体」对象来处理，就像 $\mathcal{GR}$ 是光滑有限维簇一样；也就是说，把它们等同于带某些附加结构的 $\mathcal{O}$-模层。确切地说，设 $\mathcal{GR}$ 上有 $\mathcal{O}$-模 $P$，使得 $P$ 的每个局部截面\SourcePage{197}都支在 $\mathcal{GR}$ 的某个子概形上。于是容易定义 $\operatorname{Der}\mathcal{O}_{\mathcal{GR}}$ 在 $P$ 上的连续\emph{右}作用。这样的 $P$ 带上这样的作用，就等同于 $\mathcal{GR}$ 上的 $\mathcal{D}$-模（我们还假定适当的拟凝聚性条件）。细节可见 ???。
% [Tutor Note: 原书 p.197「Details can be found in ???」中的交叉引用印作三个问号、无编号，底本如此，保留原样。]

\bdp{5.2.11} \emph{注.} 我们看到，在这个无穷维情形下，只有\emph{右} $\mathcal{D}$-模才真正作为层有意义。其原因相当有限维。事实上，设 $i:Y\hookrightarrow Z$ 是光滑流形的闭嵌入，$M$ 是 $Y$ 上的 $\mathcal{D}$-模；为把 $M$ 等同于 $i_*M$ 的子层，需要用到\emph{右} $\mathcal{D}$-模。

\bdp{5.2.12} 按 3.4.3，有群胚 $Z\operatorname{tors}_\theta(O)$，它是 $Z\operatorname{tors}_\theta(X)$ 的局部类比。选定 $\mathcal{L}\in Z\operatorname{tors}_\theta(O)$（本质上相当于选定 $\omega_O$ 的平方根）后，就定义了 $\mathcal{GR}$ 上「局部」Pfaffian 线丛 $\lambda_{\mathcal{L}}^{\operatorname{loc}}$（见 4.6）。$\mathfrak{g}\otimes K$ 经左无穷小平移在 $\mathcal{GR}$ 上的作用，提升为 2.5.1 的中心扩张 $\widetilde{\mathfrak{g}\otimes K}$ 在 $\lambda_{\mathcal{L}}^{\operatorname{loc}}$ 上的作用，使得 $\mathbf{1}\in\mathbb{C}\subset\widetilde{\mathfrak{g}\otimes K}$ 的作用为乘以 $-1$（见 4.6.5）。这给出反同态 $\bar U'\to\Gamma(\mathcal{GR},\mathcal{D}')$，其中 $\bar U'=\bar U'(\mathfrak{g}\otimes K)$ 是 2.9.4 中定义的完备化扭泛包络代数，而 $\Gamma(\mathcal{GR},\mathcal{D}')$ 是 $\mathcal{GR}$ 上 $\lambda_{\mathcal{L}}^{\operatorname{loc}}$-扭微分算子构成的环。因此，对 $\mathcal{GR}$ 上任一 $\mathcal{D}$-模 $M$，代数 $\bar U'$ 作用于 $M\lambda_{\mathcal{L}}^{-1}:=M\otimes_{\mathcal{O}_{\mathcal{GR}}}(\lambda_{\mathcal{L}}^{\operatorname{loc}})^{\otimes -1}$。故 $\Gamma(\mathcal{GR},M\lambda_{\mathcal{L}}^{-1})$ 是（左）$\bar U'$-模。

例如，考虑 $\mathcal{GR}$ 的基点处 $\delta$-函数构成的 $\mathcal{D}$-模 $I_1$。$\bar U'$-模 $\Gamma(\mathcal{GR},I_1\lambda_{\mathcal{L}}^{-1})$ 就是真空模 $Vac'$。

\bdp{5.2.13} 回忆（见 4.5.8）：$\mathcal{GR}$ 按 $G(O)$-轨道 $\operatorname{Orb}_\chi$ 分层，轨道由 $\chi\in P_+({}^{L}G)$ 标号。记 $I_\chi$ 为 $\mathcal{GR}$ 上对应于 $\overline{\operatorname{Orb}}_\chi$ 的不可约「相交上同调」$\mathcal{D}$-模。

以下是局部主定理的第一部分。

\bdp{5.2.14} \emph{定理.} $\bar U'$-模 $\Gamma(\mathcal{GR},I_\chi\lambda_{\mathcal{L}}^{-1})$ 同构于若干 $Vac'$ 的直和；并且当 $i>0$ 时 $H^i(\mathcal{GR},I_\chi\lambda_{\mathcal{L}}^{-1})=0$。

\SourcePage{198}

\emph{注.} 该定理意味着（见 5.4.8、5.4.10）：Harish-Chandra 模 $Vac'$ 是 7.8.2、7.14.1 中 Hecke 函子之 Harish-Chandra 版本的特征模。

\bdp{5.2.15} 群 $\operatorname{Aut}O$ 作用于 $\mathcal{GR}$，其李代数 $\operatorname{Der}O$ 的作用提升到 $\lambda_{\mathcal{L}}^{\operatorname{loc}}$ 上（见 4.6.7）。定理的第二部分描述 $\operatorname{Der}O$ 在 $\Gamma(\mathcal{GR},I_\chi\lambda_{\mathcal{L}}^{-1})$ 上的作用。

考虑 3.2.1 中局部 ${}^{L}\mathfrak{g}$-oper 的概形 $\mathcal{O}\mathfrak{p}_{{}^{L}\mathfrak{g}}(O)=\operatorname{Spec}A_{{}^{L}\mathfrak{g}}(O)$。把 $A_{{}^{L}\mathfrak{g}}(O)$ 简记为 $A$。与 5.2.8 一样，$\mathcal{L}$ 定义了 $\operatorname{Spec}O$ 上一族 ${}^{L}G$-oper，以 $\operatorname{Spec}A$ 为参数。这族 oper 定义了 $\operatorname{Spec}A$ 上的 ${}^{L}G$-torsor $\mathfrak{F}_A$，带有 $\operatorname{Der}O$ 的作用，且该作用与 $\operatorname{Der}O$ 在 $A$ 上的作用相容；见 3.5.4\footnote{在 3.5.4 中我们用的记号是 $\mathfrak{F}_G^0$ 而非 $\mathfrak{F}_A$，并考虑了 $\mathcal{L}$ 为 $\omega_O$ 的平方根这一「特殊」情形。}。${}^{L}G$-模 $V^\chi$ 的 $\mathfrak{F}_A$-扭是 $\operatorname{Spec}A$ 上的向量丛。记 $V^\chi_{\mathcal{L}A}$ 为其截面构成的 $A$-模；$\operatorname{Der}O$ 作用于它。

\bdp{5.2.16} \emph{定理.} 存在 $\bar U'$-模的典范同构
\begin{equation}\tag{257}
\Gamma(\mathcal{GR},I_\chi\lambda_{\mathcal{L}}^{-1})\xrightarrow{\sim}Vac'\otimes_AV^\chi_{\mathcal{L}A},
\end{equation}
它与 $\operatorname{Der}O$ 的作用相容。

这里我们用到 $Vac'$ 上的 $A$-模结构，它来自 Feigin--Frenkel 同构 (80)。

\bdp{5.2.17} 关于证明略说几句。整体定理由局部定理经一个容易的局部--整体论证得到，该论证与 2.8 中所用者类似。局部定理的证明基于下述两个关键结构之间的相互作用：
\begin{enumerate}[label=(\roman*),leftmargin=2.2em]
\item \emph{Satake 等价}（[Gi95]、[MV]）：${}^{L}G$ 的表示张量范畴，与由诸 $I_\chi$ 生成、带「卷积」张量结构的 $\mathcal{GR}$ 上 $\mathcal{D}$-模范畴之间的等价。
\item \emph{「重整化」包络代数} $U^\natural$。代数态射 $\bar U'\to\Gamma(\mathcal{GR},\mathcal{D}')$ 既非单射（它消去\SourcePage{199}$Vac'$ 在 $\bar U'$ 的中心 $\mathfrak{z}$ 中的零化子 $I$），也非满射（其像不含 $\operatorname{Der}O$）。我们把它分解为 $\bar U'\to U^\natural\to\Gamma(\mathcal{GR},\mathcal{D}')$，其中 $U^\natural$ 是在 $\bar U'/I\bar U'$ 上「添加」3.6.5 的代数胚 $I/I^2$ 而得（$\mathfrak{z}_{\mathfrak{g}}(O)=\mathfrak{z}/I\subset\bar U'/I\bar U'$ 与 $I/I^2$ 之间的交换关系来自 $I/I^2$ 上的代数胚结构，它们几乎是 Heisenberg 型的）。真空表示 $Vac'$ 作为 $U^\natural$-模不可约；$\Gamma(\mathcal{GR},I_\chi\lambda^{-1})$ 亦然，$\chi\in P_+({}^{L}G)$。
\end{enumerate}

\bdp{5.2.18} 以下是 5.2.16 的证明思路（我们假定 5.2.14）。置 $\mathfrak{z}:=\mathfrak{z}_{\mathfrak{g}}(O)$。考虑 $\mathfrak{z}$-模 $V^\chi_{\mathcal{L}\mathfrak{z}}:=\operatorname{Hom}_{\bar U'}(Vac',\allowbreak\Gamma(\mathcal{GR},I_\chi\lambda_{\mathcal{L}}^{-1}))$，于是 $\Gamma(\mathcal{GR},I_\chi\lambda_{\mathcal{L}}^{-1})=Vac'\otimes_{\mathfrak{z}}V^\chi_{\mathcal{L}\mathfrak{z}}$。某种 Tannakian 形式体系连同 Satake 等价，给出 $\operatorname{Spec}\mathfrak{z}$ 上典范的 ${}^{L}G$-torsor $\mathfrak{F}_{\mathfrak{z}}$，使得诸 $V^\chi_{\mathcal{L}\mathfrak{z}}$ 是 $V^\chi$ 的 $\mathfrak{F}$-扭。$\Gamma(\mathcal{GR},I_\chi\lambda_{\mathcal{L}}^{-1})$ 上的 $U^\natural$-模结构，定义了李代数胚 $I/I^2$ 在 $\mathfrak{F}_{\mathfrak{z}}$ 上的作用。一些额外的几何考虑在 $\mathfrak{F}_{\mathfrak{z}}$ 上定义了典范的 $B$-结构，它关于 $\operatorname{Der}O\subset I/I^2$ 的作用满足「oper」性质。于是 3.5、3.6 的结果给出典范等同 $(\operatorname{Spec}\mathfrak{z},\mathfrak{F}_{\mathfrak{z}})\simeq(\operatorname{Spec}A,\mathfrak{F}_A)$，其中 $A\simeq\mathfrak{z}$ 是 Feigin--Frenkel 同构。证毕。

\bdp{5.2.19} 我们需要它吗？？？
% [Tutor Note: 原书 p.199 此处标题为「DO WE NEED IT???」，是作者留在正文中的自问（三个问号照录），底本如此，保留原样。]

下面给出 $M$ 的一个直接构造，不用扭 $\mathcal{D}$-模。对 $x\in X$，考虑概形 $\operatorname{Bun}_{G,\bar x}$（见 2.3.1）。对 $\mathcal{L}\in Z\operatorname{tors}_\theta(X)$，记 $\lambda_{\mathcal{L},\bar x}$ 为线丛 $\lambda_{\mathcal{L}}$ 到 $\operatorname{Bun}_{G,\bar x}$ 的拉回。设 $\widetilde{\mathfrak{g}\otimes K_x}$ 是 2.5.1 中 $\mathfrak{g}\otimes K_x$ 的中心扩张。于是 $\mathfrak{g}\otimes K_x$ 在 $\operatorname{Bun}_{G,\bar x}$ 上的作用，典范地提升为 $\widetilde{\mathfrak{g}\otimes K_x}$ 在 $\lambda_{\mathcal{L},\bar x}$ 上的作用，其中 $\mathbf{1}\in\mathbb{C}$ 的作用为恒等（见 4.4.12）。记 $\operatorname{Bun}_{G,\mathcal{L},\bar x}$ 为 $\operatorname{Bun}_{G,\bar x}$ 上对应于 $\lambda_{\mathcal{L},\bar x}$ 的 $\mathbb{G}_m$-torsor 的全空间。我们有 Harish-Chandra 对 $(\widetilde{\mathfrak{g}\otimes K_x},\mathbb{G}_m\times G(O_x))$，其中 $\operatorname{Lie}\mathbb{G}_m=\mathbb{C}\subset\widetilde{\mathfrak{g}\otimes K_x}$。$\widetilde{\mathfrak{g}\otimes K_x}$ 在 $\operatorname{Bun}_{G,\mathcal{L},\bar x}$ 上的作用，按显然的方式延拓为该对的作用。

注意 $\operatorname{Bun}_G=\mathbb{G}_m\times G(O_x)\backslash\operatorname{Bun}_{G,\mathcal{L},\bar x}$。因此按 1.2.4 与 1.2.6，有函子 $\Delta_{\mathcal{L}}:\allowbreak\left(\widetilde{\mathfrak{g}\otimes K_x},\mathbb{G}_m\times G(O_x)\right)\allowbreak\operatorname{mod}\to\mathcal{M}^\ell(\operatorname{Bun}_G)$。

\SourcePage{200}

把投影 $\mathbb{G}_m\times G(O_x)\to\mathbb{G}_m$ 看作一个特征标；记 $Vac^\sim$ 为相应的诱导 Harish-Chandra 模。有
\begin{equation}\tag{258}
M_{\mathcal{L}}=\Delta_{\mathcal{L}}(Vac^\sim).
\end{equation}

下面确定 $M_{\mathcal{L}}$ 上的 $A_{{}^{L}\mathfrak{g}}(X)$-模结构。$\operatorname{End}(Vac^\sim)=\mathfrak{z}_{\mathfrak{g}}(O_x)$ 在 $\Delta_{\mathcal{L}}(Vac^\sim)$ 上的作用，经 Feigin--Frenkel 同构 $\varphi_{O_x}$（见 3.2.2），等同于一个 $A_{{}^{L}\mathfrak{g}}(O_x)$-作用。该作用经 $A_{{}^{L}\mathfrak{g}}(X)$ 这个商分解。

\bdsec{5.3. Satake 等价}

本节回顾基本事实与构造，并固定记号。细节与证明见 [MV]。[MV] 的作者用反常层；我们用 $\mathcal{D}$-模。

\bdp{5.3.1} 考虑仿射（或环路）Grassmann 流形 $\mathcal{GR}=G(K)/G(O)$（按惯例 $K=\mathbb{C}((t))$，$O=\mathbb{C}[[t]]$）；它是形式光滑、归纳射影的归纳概形（见 4.5.1）。它带有按 $G(O)$-轨道 $\operatorname{Orb}_\chi$（$\chi\in P_+({}^{L}G)$）的分层（见 4.5.8）。每个分层都是 $\operatorname{Aut}^0O$-不变的。

在 4.5.10 中我们引入了 $\mathcal{GR}$ 连通分支的奇偶性概念。按 4.5.11，
\begin{equation}\tag{259}
\begin{array}{c}
\text{$\mathcal{GR}$ 的偶（相应地奇）连通分支的所有分层，}\\[2pt]
\text{维数都是偶（相应地奇）的。}
\end{array}
\end{equation}

\bdp{5.3.2} \emph{引理.}
\begin{enumerate}[label=(\roman*),leftmargin=2.2em]
\item 每个分层 $\operatorname{Orb}_\chi$ 连通且单连通。
\item $\operatorname{Orb}_\chi$ 上任一光滑 $\mathcal{D}$-模是常值的。
\item $\operatorname{Orb}_\chi$ 只在偶次有上同调。
\end{enumerate}

\emph{证明.} 记 $\operatorname{Stab}_x$ 为 $x\in\mathcal{GR}$ 在 $G(O)$ 中的稳定化子。$\operatorname{Stab}_x$ 在 $G(O/tO)=G$ 中的像是抛物子群 $P_x$，而态射 $G(O)/\operatorname{Stab}_x\to G/P_x$ 是局部平凡纤维化，其纤维同构于仿射空间。于是 (i) 与 (iii) 显然。注意 $\overline{\operatorname{Orb}}_\chi$ 射影，且按 (259)，$\overline{\operatorname{Orb}}_\chi\setminus\operatorname{Orb}_\chi$ 的余维 $\ge 2$。故按
% 本单元止于 p.200 末句「So by」；下文（the projection formula …）见 p.201，由下一单元承接。

% ===== units/ch05_u02.tex =====
% ch05_u02 —— 原书 p.201–p.209（5.3.3 的后半句 – 5.3.17）
% 首句承接 ch05_u01 末尾的「故按」（原书 p.200 末 "So by"），拼接后读作「故按 Deligne 定理，……」。
\SourcePage{201}

Deligne 定理\footnote{不用 Deligne 定理也可注意到：$\operatorname{Orb}_\chi$ 上任一向量丛的解析截面都是代数的。把它用于带可积联络的向量丛在 $\operatorname{Orb}_\chi$ 上的水平解析截面，即知 (ii) 由 (i) 推出。}，$\operatorname{Orb}_\chi$ 上的光滑 $\mathcal{D}$-模具有正则奇点，因此 (ii) 由 (i) 推出。\hfill$\Box$

记 $\mathcal{P}$ 为 $\mathcal{GR}$ 上沿我们的分层光滑的凝聚（等价地，完整）$\mathcal{D}$-模构成的范畴。

\bdp{5.3.3} \emph{命题.}

\begin{enumerate}[label=(\roman*),leftmargin=2.2em]
\item 范畴 $\mathcal{P}$ 是半单的。
\item 若 $M\in\mathcal{P}$ 支撑在偶（相应地，奇）连通分支上，则当 $a$ 为奇（相应地，偶）数时 $H^{a}_{DR}(\mathcal{GR},M)=0$。
\end{enumerate}

\emph{证明.} 记 $I_\chi$ 为 $\operatorname{Orb}_\chi$ 上 $\mathbb{C}$-向量空间构成的相交上同调反常层。记 $\mathcal{GR}_{(\chi)}$ 为 $\mathcal{GR}$ 中含 $\operatorname{Orb}_\chi$ 的连通分支，$p(\chi)$ 为 $\mathcal{GR}_{(\chi)}$ 的奇偶性。由 Lusztig（[Lu83] 的定理 11c），$I_\chi$ 具有下述性质：上同调层 $H^{i}(I_\chi)$ 为零，除非 $i\bmod 2=p(\chi)$。记 $C$ 为 $D^{b}(\mathcal{GR}_{(\chi)})$ 中具有该性质、且沿我们的分层光滑的全部对象构成的范畴。由 (259) 与 5.3.2(iii)，对任意 $M,N\in C$，只要 $i\bmod 2\neq p(\chi)$ 就有 $H^{i}(\mathcal{GR}_{(\chi)},M)=0$；且对奇数 $i$ 有 $\operatorname{Ext}^{i}(M,N^{*})=0$（这里 $N^{*}$ 是 $N$ 的 Verdier 对偶）。特别地，只要 $i\bmod 2\neq p(\chi)$ 就有 $H^{i}(\mathcal{GR},I_\chi)=0$；又 $\operatorname{Ext}^{1}(I_{\chi_1},I_{\chi_2})=0$。用 5.3.2(ii) 即得命题。\hfill$\Box$

\bdp{5.3.4} 由 5.3.2(ii)，$\mathcal{P}$ 的单对象是各层 $\operatorname{Orb}_\chi$ 的「相交上同调」$\mathcal{D}$-模 $I_\chi$。于是 5.3.3(i) 表明，$\mathcal{P}$ 的任一对象都具有 $G(O)$-等变或 $\operatorname{Aut}^{0}O\ltimes G(O)$-等变的 $\mathcal{D}$-模结构。这种结构唯一，且任何态射都与它相容（因为我们的群连通）。我们看到，\SourcePage{202}
$\mathcal{P}$ 等同于 $\mathcal{GR}$ 上 $G(O)$-等变或 $\operatorname{Aut}^{0}O\ltimes G(O)$-等变的凝聚 $\mathcal{D}$-模构成的范畴。

\emph{注.} $G(O)$-等变结构的存在性也可直接由下述事实推出：$G(O)$ 连通，且 $\operatorname{Hom}(G(O),\mathbb{G}_m)=0$（以及 5.3.2(ii)）。此时无需援引 5.3.3(i)，因而也无需援引 Lusztig 定理（那是一个深刻的结果）。

\bdp{5.3.5} 范畴 $\mathcal{P}$ 带有典范的张量结构。有两种方式描述它：「卷积」构造（见 5.3.5 -- 5.3.9）与「融合」构造（先作 5.3.10 -- 5.3.12 的若干准备，再在 5.3.13 -- 5.3.16 中给出）；两种定义的等价性见 5.3.17。我们从卷积图景\footnote{以下内容是 Ginzburg 的拓扑构造 [Gi95] 的代数版本；辨认这两种构造的工作留给感兴趣的读者。}开始。我们需要定义卷积积函子 $\circledast:\mathcal{P}\times\mathcal{P}\to\mathcal{P}$、$\circledast$ 的结合约束，以及交换约束。

由 [MV]，函子 $\circledast$ 定义如下。记 $G(K)\times_{G(O)}\mathcal{GR}$ 为 $G(K)\times\mathcal{GR}$ 对 $G(O)$ 的商，其中 $u\in G(O)$ 通过 $(g,x)\mapsto(gu^{-1},ux)$ 作用在 $G(K)\times\mathcal{GR}$ 上。态射 $p:G(K)\times_{G(O)}\mathcal{GR}\to G(K)/G(O)=\mathcal{GR}$，$(g,x)\mapsto g\bmod G(O)$，是纤维为 $\mathcal{GR}$ 的局部平凡纤维化；它附属于主 $G(O)$-丛 $G(K)\to\mathcal{GR}$ 与 $G(O)$ 在 $\mathcal{GR}$ 上的作用。故 $G(K)\times_{G(O)}\mathcal{GR}$ 是 $\mathcal{GR}\times\mathcal{GR}$ 的一个扭形式。设 $M,N\in\mathcal{P}$。用 $M$ 上的 $G(O)$-等变结构，可定义 $G(K)\times_{G(O)}\mathcal{GR}$ 上的一个 $\mathcal{D}$-模 $M\boxtimes'N$，它是 $M\boxtimes N$ 的「扭形式」。于是

\begin{equation}\tag{260}
M\circledast N=m_{*}(M\boxtimes'N),
\end{equation}

其中 $m:G(K)\times_{G(O)}\mathcal{GR}\to\mathcal{GR}$ 来自 $G(K)$ 在 $\mathcal{GR}$ 上的作用。

\bdp{5.3.6} \emph{奇迹定理（Lusztig）.} 若 $M,N\in\mathcal{P}$，则 $M\circledast N\in\mathcal{P}$。\hfill$\Box$

\SourcePage{203}

\emph{注.} 非平凡的论断是 $M\circledast N$ 为 $\mathcal{D}$-模（而不只是导出范畴的对象）。由于该 $\mathcal{D}$-模凝聚且 $G(O)$-等变，它属于 $\mathcal{P}$。

\emph{注.} G. Lusztig 对 5.3.6 的原证明 [Lu83] 基于显式估计，说明 $m$ 限制到 $\mathcal{GR}$ 中两个层的（扭）积上是半小的。很久以后 D. Gaitsgory 指出 [Ga2001]：只要把融合张量积定义中涉及的几何结构考虑进来，5.3.6 便立即成立。

于是我们定义了 $\circledast:\mathcal{P}\times\mathcal{P}\to\mathcal{P}$。$\circledast$ 的结合约束按显然的方式定义。交换约束将在 5.3.8 中定义。

\bdp{5.3.7} \emph{注.} (i) 设 $G(K)$ 换成一个归纳仿射的群归纳概形 $\mathcal{G}$，$G(O)$ 换成其闭子群概形 $\mathcal{K}$；并设 $\mathcal{G}/\mathcal{K}$ 是归纳有限型的归纳概形。5.3.5 中 $\circledast:\mathcal{P}\times\mathcal{P}\to\mathcal{P}$ 的构造基于奇迹 5.3.6。一般地，$\mathcal{G}/\mathcal{K}$ 上 $\mathcal{K}$-等变 $\mathcal{D}$-模的范畴上没有卷积；必须考虑某个导出范畴 $\mathcal{H}$（Hecke 幺半范畴；见 7.6.1 与 7.11.17）。这是一个带 t-结构的三角范畴，其核心是 $\mathcal{G}/\mathcal{K}$ 上 $\mathcal{K}$-等变 $\mathcal{D}$-模的范畴；一般地，$\circledast:\mathcal{H}\times\mathcal{H}\to\mathcal{H}$ 不是 t-正合的，且 $\circledast$ 没有交换约束。在 $(G(K),G(O))$ 的情形，由 5.3.6，函子 $\circledast$ 是 t-正合的，而 $\mathcal{H}$ 的核心是 $\mathcal{P}$ 的归纳对象构成的范畴。

(ii) 上述 $\mathcal{H}$ 的构造是 §7 中发展的「Hecke 模式」的一部分。后面我们将看到，即使在 $(G(K),G(O))$ 这一奇迹般好的情形，该模式也是有用的（或许不可或缺）。

\bdp{5.3.8} 我们来定义 $\circledast$ 的交换约束。设 $\theta:G\to G$ 是把每个支配权映到其对偶的自同构。$G$ 的反自同构 $\theta'(g):=\theta(g)^{-1}$ 给出 $\mathcal{H}$ 的反自同构 $\theta'_{\mathcal{H}}$\SourcePage{204}，故对任意 $M,N\in\mathcal{H}$ 有典范同构

\[
l_{M,N}:\theta'_{\mathcal{H}}(M\circledast N)\xrightarrow{\sim}\theta'_{\mathcal{H}}(N)\circledast\theta'_{\mathcal{H}}(M).
\]

对任意 $M\in\mathcal{P}\subset\mathcal{H}$ 有典范同构 $e_M:M\xrightarrow{\sim}\theta'_{\mathcal{H}}(M)$。按 5.3.3(i)，为定义 $e_M$ 只需考虑 $M=I_\chi$ 的情形。$\theta'$ 在 $G(K)$ 上的作用保持由双 $G(O)$-类给出的分层 $G(K)_\chi$（这里 $G(K)_\chi$ 是 $\operatorname{Orb}_\chi\subset G(K)/G(O)$ 的原像）。于是有诱导的自同构 $\theta'_\chi$。作为 $\mathcal{H}$ 的对象，我们的 $I_\chi$ 是 $G(K)$ 上的 $\Omega$-复形 $\Omega_{G(K)_\chi}[\dim\operatorname{Orb}_\chi]$。现在 $e_{I_\chi}$ 是 $\theta'_\chi$ 在 $\Omega_{G(K)_\chi}$ 上的作用。

对 $M,N\in\mathcal{P}$ 定义

\begin{equation}\tag{261}
s:M\circledast N\xrightarrow{\sim}N\circledast M
\end{equation}

为如下复合：

\[
M\circledast N\xrightarrow{\sim}\theta'_{\mathcal{H}}(M\circledast N)\xrightarrow{\sim}\theta'_{\mathcal{H}}(N)\circledast\theta'_{\mathcal{H}}(M)\xrightarrow{\sim}N\circledast M,
\]

其中第一个箭头是对应于 $M\circledast N$ 的同构 $e$，其余箭头是 $l_{M,N}$ 与 $e_N^{-1}\circledast e_M^{-1}$。

\bdp{5.3.9} \emph{命题.} $s$ 是卷积张量积 $\circledast$ 的交换约束。

\emph{证明.} 在下文 5.3.17 中，我们把卷积张量积与融合张量积等同起来，且该等同与所有约束相容。由于后一组数据显然在 $\mathcal{P}$ 上定义了张量范畴结构，证毕。\hfill$\Box$

于是我们定义了 $\mathcal{P}$ 上预期的卷积张量结构。

\bdp{5.3.10} $\mathcal{P}$ 上张量结构的融合描述\footnote{该构造表面上涉及一条曲线 $X$，但实际上纯粹是局部的。}基于 4.3.14 中「空间」$\operatorname{GRAS}=\operatorname{GRAS}_G$ 上重要的手性半群结构。该结构可描述如下。

\SourcePage{205}

(i) 对 $\mathbb{C}$-代数 $R$ 与 $S\in\Sigma(R)$（我们沿用 4.3.11 的记号，故 $S$ 是 $X\otimes R$ 的一个在 $\operatorname{Spec}R$ 上有限且平坦的子概形），有子集 $\operatorname{GRAS}(R)_S\subset\operatorname{GRAS}(R)$，它由对 $(\mathcal{F},\gamma)$ 构成，其中 $\mathcal{F}$ 是 $X\otimes R$ 上的 $G$-丛，$\gamma$ 是 $\mathcal{F}$ 在 $S$ 的补集上的截面。

(ii) 若 $S$ 是子概形 $S_i$（$i\in I$）的不交并，则有典范等同

\begin{equation}\tag{262}
\operatorname{GRAS}(R)_S\xrightarrow{\sim}\prod_i\operatorname{GRAS}(R)_{S_i}.
\end{equation}

确切地说，我们把 $(\mathcal{F},\gamma)$ 与族 $(\mathcal{F}_i,\gamma_i)$（$i\in I$）等同，其中 $(\mathcal{F}_i,\gamma_i)\in\operatorname{GRAS}(R)_{S_i}$ 在 $S_{i'}$（$i'\neq i$）之并的补集上与 $(\mathcal{F},\gamma)$ 一致。

数据 (i)、(ii) 具有如下性质：

a. 若对 $S_1,S_2\in\Sigma(R)$ 有 $S_{1\mathrm{red}}\subset S_{2\mathrm{red}}$，则 $\operatorname{GRAS}(R)_{S_1}\subset\operatorname{GRAS}(R)_{S_2}$。$\operatorname{GRAS}(R)_S$（$S\in\Sigma(R)$）之并等于 $\operatorname{GRAS}(R)$。故 $\operatorname{GRAS}(R)_S$ 构成 $\operatorname{GRAS}(R)$ 上的滤链。该滤链具有函子性（关于 $R$）。

b. 同构 (ii) 也具有函子性，并按显然的方式与 $I$ 的细分相容。

c. 由

\[
\mathcal{GR}_\Sigma(R):=\{(S,\mathcal{F},\gamma)\mid S\in\Sigma(R),\ (\mathcal{F},\gamma)\in\operatorname{GRAS}(R)_S\}
\]

定义的子函子 $\mathcal{GR}_\Sigma\subset\Sigma\times\operatorname{GRAS}$ 是 $\Sigma$ 上形式光滑的归纳概形。

\emph{注.} 我们来解释对任意仿射代数群 $G$，$\mathcal{GR}_\Sigma=\mathcal{GR}^{G}_\Sigma$ 为何是归纳概形。此外我们将证明：$\mathcal{GR}_\Sigma$ 是归纳有限型的；且当 $G$ 约化时 $\mathcal{GR}_\Sigma$ 归纳固有。先考虑 $G=GL_n$ 的情形。此时 $\mathcal{GR}_\Sigma$ 是诸 $\mathcal{GR}_{\Sigma,k}$ 的正向极限，其中 $\mathcal{GR}_{\Sigma,k}$ 参数化由有限子概形 $D\subset X$ 与子层 $\mathcal{E}\subset\mathcal{O}_X^{n}(kD)$ 构成、且满足 $\mathcal{E}\supset\mathcal{O}_X^{n}(-kD)$ 的对。态射 $\mathcal{GR}_{\Sigma,k}\to\Sigma$ 固有，故 $\mathcal{GR}_\Sigma$ 归纳固有。如定理 4.5.1 的证明中所说明，为把一般情形化到\SourcePage{206} $GL_n$ 的情形，只需证明：若 $G\subset G'$ 且 $G'/G$ 仿射（相应地，拟仿射），则态射 $\mathcal{GR}^{G}_\Sigma\to\mathcal{GR}^{G'}_\Sigma$ 是闭（相应地，局部闭）嵌入。这很容易。

\bdp{5.3.11} 对有限集 $J$，有态射 $X^{J}\to\Sigma$，它把 $(x_j)\in X^{J}$ 对应到与除子 $\sum_j x_j$ 相应的子概形 $D\subset X$。记 $\mathcal{GR}_{X^{J}}$ 为 $\mathcal{GR}_\Sigma$ 与 $X^{J}$ 在 $\Sigma$ 上的纤维积。于是 $\mathcal{GR}_{X^{J}}$ 的一个 $R$-点是一个族 $((x_j),\mathcal{F},\gamma)$，其中 $(x_j)\in X^{J}(R)$，$\mathcal{F}$ 是 $X\otimes R$ 上的 $G$-丛，$\gamma$ 是 $\mathcal{F}$ 在诸 $x_j$ 的图像之并的补集上的截面。我们的 $\mathcal{GR}_{X^{J}}$ 是 $X^{J}$ 上形式光滑、归纳固有的归纳概形（见 5.3.10 末的注）。

由 4.5.2，$\mathcal{GR}_X$ 在 $x\in X(\mathbb{C})$ 上的纤维与归纳概形 $\mathcal{GR}_x:=G(K_x)/G(O_x)$ 之间有典范同构。故由 5.3.10(ii)，$\mathcal{GR}_{X^{J}}$ 在 $(x_j)\in X^{J}(\mathbb{C})$ 上的纤维等于 $\prod_{x\in S}\mathcal{GR}_x$，其中 $S$ 是子集 $\{x_j\}\subset X$。

下面的 $\mathcal{GR}_X$ 的描述将有用。考虑 $X$ 上「形式参数」的概形 $X^{\wedge}$（其点是光滑态射 $\operatorname{Spec}O\to X$，见 2.6.5）。这是 $X$ 上的 $\operatorname{Aut}^{0}O$-torsor；坐标的选择，即某个开集 $U\subset X$ 上平展的 $\mathbb{A}^{1}$-值映射，给出 $X^{\wedge}$ 在 $U$ 上的平凡化。现在 $\mathcal{GR}_X$ 是 $\mathcal{GR}$ 的 $X^{\wedge}$-扭（关于 $\operatorname{Aut}^{0}O$ 在 $\mathcal{GR}$ 上的作用）。$\mathcal{GR}$ 的分层给出 $\mathcal{GR}_X$ 的一个分层，其层为 $X$ 上光滑的 $\operatorname{Orb}_{\chi X}$。

\bdp{5.3.12} 为便于将来引用，我们列出 $\mathcal{GR}_{X^{J}}$ 之间的一些相容性；它们直接由 5.3.10 推出。

a. 对满射 $\pi:J\twoheadrightarrow J'$，有显然的笛卡尔图表

\begin{equation}\tag{263}
\begin{tikzcd}
\mathcal{GR}_{X^{J'}} \arrow[r,"\tilde\Delta^{(\pi)}"] \arrow[d] & \mathcal{GR}_{X^{J}} \arrow[d]\\
X^{J'} \arrow[r,"\Delta^{(\pi)}"] & X^{J}
\end{tikzcd}
\end{equation}

\SourcePage{207}

其中 $\Delta^{(\pi)}$ 是 $\pi$-对角嵌入。若 $|J'|=1$，则有 $\Delta^{(J)}:X\hookrightarrow X^{J}$ 与 $\tilde\Delta^{(J)}:\mathcal{GR}_X\hookrightarrow\mathcal{GR}_{X^{J}}$。

b. 设 $\nu^{(J)}:U^{(J)}\hookrightarrow X^{J}$ 是对角除子的补集。由 5.3.10(ii)，$X^{J}$-归纳概形 $\mathcal{GR}_{X^{J}}$ 与 $(\mathcal{GR}_X)^{J}$ 在 $U^{(J)}$ 上的限制有典范等同。因此有笛卡尔图表

\begin{equation}\tag{264}
\begin{tikzcd}
(\mathcal{GR}_X)^{J}|_{U^{(J)}} \arrow[r,"\tilde\nu^{(J)}"] \arrow[d] & \mathcal{GR}_{X^{J}} \arrow[d]\\
U^{(J)} \arrow[r,"\nu^{(J)}"] & X^{J}
\end{tikzcd}
\end{equation}

\bdp{5.3.13} 现在我们准备定义 $\mathcal{P}$ 上的融合张量结构。这归结为对任意有限非空集 $J$ 构造张量积函子\footnote{这里 $\mathcal{P}^{\otimes J}$ 表示 $\mathcal{P}$ 的 $J$ 个副本的张量积（由于 $\mathcal{P}$ 半单，张量积的定义是清楚的）。}

\begin{equation}\tag{265}
\circledast_J:\mathcal{P}^{\otimes J}\to\mathcal{P}
\end{equation}

并给出等同

\begin{equation}\tag{266}
\circledast_J=\circledast_{J'}\Bigl(\bigotimes_{j'\in J'}\bigl(\circledast_{\pi^{-1}(j')}\bigr)\Bigr)
\end{equation}

对任意满射 $J\xrightarrow{\pi}J'$。

构造如下。

\bdp{5.3.14} 由于任一 $M\in\mathcal{P}$ 是 $\operatorname{Aut}^{0}O$-等变的，它定义了 $\mathcal{GR}_X$ 上的一个 $\mathcal{D}$-模（见 5.3.11 末对 $\mathcal{GR}_X$ 的描述）。记 $M_X\in D(\mathcal{GR}_X)$（$:=DM(\mathcal{GR}_X)$）为它在导出范畴中的平移 1。换言之，对上述任一开集 $U$ 与 $X^{\wedge}$ 在 $U$ 上的平凡化 $\theta$，有 $M_U=\pi_\theta^{!}M$；其中 $M_U:=M_X|_{\mathcal{GR}_U}$，$\pi_\theta:\mathcal{GR}_U\to\mathcal{GR}$ 是对应于 $\theta$ 的投影。我们用 $M$ 上的 $\operatorname{Aut}^{0}O$-作用把这些对象粘合起来。函子 $\mathcal{P}\to D(\mathcal{GR}_X)$，$M\mapsto M_X$，是全忠实的。其本质像由（平移 1 的）$\mathcal{D}$-模构成，这些 $\mathcal{D}$-模同构于「相交上同调」$\mathcal{D}$-模 $I_{\chi X}$ 的（有限多个）副本的直和；$I_{\chi X}$ 对应于 $\operatorname{Orb}_{\chi X}$ 上的平凡局部系统。

\SourcePage{208}

设 $\{M_j\}_{j\in J}$ 是 $\mathcal{P}$ 中对象构成的族。用 (264)，把 $\boxtimes M_{jX}|_{U^{(J)}}$ 解释为 $\mathcal{GR}_{X^{J}}|_{U^{(J)}}$ 上的、平移 $|J|$ 的 $\mathcal{D}$-模。记 $\boxtimes M_{jX}\in D(\mathcal{GR}_{X^{J}})$ 为它到 $\mathcal{GR}_{X^{J}}$ 的极小（即 $\tilde\nu^{(J)}_{!*}$-）扩张。这是 $\mathcal{GR}_{X^{J}}$ 上平移 $|J|$ 的 $\mathcal{D}$-模。因此我们定义了函子

\begin{equation}\tag{267}
\boxtimes_J:\mathcal{P}^{\otimes J}\to D(\mathcal{GR}_{X^{J}}),\qquad \otimes M_j\mapsto\boxtimes M_{jX},
\end{equation}

它显然是全忠实的。

\bdp{5.3.15} \emph{命题.}（[MV]）

对任意 $\pi:J\twoheadrightarrow J'$，复形 $\tilde\Delta^{(\pi)!}(\boxtimes M_{jX})\in D(\mathcal{GR}_X^{(J')})$ 属于 $\boxtimes_{J'}$ 的本质像。\hfill$\Box$
% [Tutor Note: 原书 p.208 此处印作 $D(\mathcal{GR}_X^{(J')})$，与 (267) 中的 $\mathcal{GR}_{X^{J}}$ 写法不一致；底本如此，保留原样。]

\bdp{5.3.16} 我们得到函子

\begin{equation}\tag{268}
\circledast_\pi:\mathcal{P}^{\otimes J}\to\mathcal{P}^{\otimes J'},
\end{equation}

使得 $\boxtimes_{J'}\circledast_\pi=\tilde\Delta^{(\pi)!}\boxtimes_J$。特别地，当 $|J'|=1$ 时得到函子 $\circledast_J:\mathcal{P}^{\otimes J}\to\mathcal{P}$，它就是我们的张量积函子 (265)。显然的等同 $\circledast_\pi=\bigotimes_{j'\in J'}\bigl(\circledast_{\pi^{-1}(j')}\bigr)$（看我们在 $U^{(J')}$ 上的 $\mathcal{D}$-模）与标准同构 $\Delta^{(J)!}=(\Delta^{(\pi)}\Delta^{(J')})^{!}=\Delta^{(J')!}\Delta^{(\pi)!}$ 给出相容性同构 (266)。故 $\mathcal{P}$ 是张量范畴。容易看出 $I_0$ 是 $\mathcal{P}$ 中的单位对象。

\bdp{5.3.17} 我们来把 $\mathcal{P}$ 上的卷积与融合张量结构等同起来。下面在本小节中，用 $\circledast^{c}$ 表示卷积张量积，用 $\circledast^{f}$ 表示 $\mathcal{P}$ 上的融合张量积。我们要对 $M,N\in\mathcal{P}$ 构造与结合约束、交换约束相容的典范同构 $M\circledast^{c}N\xrightarrow{\sim}M\circledast^{f}N$。\footnote{该构造借自 [MV]，那里写得更详细；不过那里没有考虑交换约束 5.3.8。}

设 $\mathcal{GR}'_{X^{2}}$ 是 $X^{2}$ 上的归纳概形，使得 $\mathcal{GR}'_{X^{2}}(R)$ 是族 $(x_1,x_2,\mathcal{F}_1,\mathcal{F}_2,\gamma_1,\gamma_2)$ 的集合，其中 $x_1,x_2\in X(R)$，$\mathcal{F}_1,\mathcal{F}_2$ 是 $X\otimes R$ 上的 $G$-torsor，$\gamma_1$ 是 $\mathcal{F}_1$ 在\SourcePage{209} $x_1$ 的图像之补集上的截面，$\gamma_2$ 是 $\mathcal{F}_1\to\mathcal{F}_2$ 在 $x_2$ 的图像之补集上的同构。有投影 $q:\mathcal{GR}'_{X^{2}}\to\mathcal{GR}_{X^{2}}$，它把上述数据映到 $((x_1,x_2),\mathcal{F}_2,\gamma_2\gamma_1)$。该投影是归纳固有的；在 $U:=X^{2}\setminus\{\text{对角线}\}$ 上它是同构。\footnote{在对角线上，$q$ 的纤维同构于 $\mathcal{GR}$；更确切地说，由 $(x_1,x_2,\mathcal{F}_1,\mathcal{F}_2,\gamma_1,\gamma_2)\mapsto(x_1,x_2,\mathcal{F}_1,\gamma_1,\mathcal{F}_2,\gamma_2\gamma_1)$ 定义的闭嵌入 $\mathcal{GR}'_{X^{2}}\to(\mathcal{GR}_X)\times_X(\mathcal{GR}_{X^{2}})$ 限制到对角线 $X\hookrightarrow X^{2}$ 上成为同构。故 $q$ 为同构的最大开子集形如 $\mathcal{GR}_{X^{2}}\setminus Z$，其中 $Z$ 的余维为 1；这是一个无穷维现象。}
% [Tutor Note: 原书 p.209 脚注中闭嵌入的目标印作 $(\mathcal{GR}_X)\times_X(\mathcal{GR}_{X^2})$，按上下文疑为 $(\mathcal{GR}_X)\times_X(\mathcal{GR}_X)$ 之误；底本如此，保留原样。]

记 $M_X\boxtimes'N_X\in D(\mathcal{GR}'_{X^{2}})$ 为 $M_X\boxtimes N_X|_U$ 到 $\mathcal{GR}'_{X^{2}}$ 的极小扩张。这是 $\mathcal{GR}'_{X^{2}}$ 上平移 2 的 $\mathcal{D}$-模。由 [MV]，$U$ 上的显然等同（唯一地）延拓为典范同构

\begin{equation}\tag{269}
q_{*}(M_X\boxtimes'N_X)\xrightarrow{\sim}M_X\boxtimes N_X.
\end{equation}

现在 $\mathcal{GR}'_{X^{2}}$ 是 $(\mathcal{GR}_X)^{2}$ 的扭形式。事实上，$\mathcal{F}_1$ 在 $x_2$ 的形式邻域上的平凡化给出上述数据 $(\mathcal{F}_2,\gamma_2)$ 与 $\mathcal{GR}_x$ 的一个等同。这些平凡化连同 $x_2$ 处的形式参数构成 $\mathcal{GR}_X\times X$ 上的 $\operatorname{Aut}^{0}O\ltimes G(O)$-torsor，而 $\mathcal{GR}'_{X^{2}}$ 与 $\mathcal{GR}$ 的相应扭等同。故 $M_X\boxtimes'N_X$ 是 $M_X\boxtimes N$ 的「扭形式」。把这一图景限制到对角线 $X\hookrightarrow X\times X$，我们看到：$q:\mathcal{GR}'_{X^{2}}\to\mathcal{GR}_{X^{2}}$ 到 $X$ 的拉回与 (260) 中态射 $m:G(K)\times_{G(O)}\mathcal{GR}\to\mathcal{GR}$ 的 $X^{\wedge}$-扭重合；而 $M_X\boxtimes'N_X$ 到 $X$ 在 $\mathcal{GR}'_{X^{2}}$ 中前像的拉回等于 $(M\boxtimes'N)_X$，其中 $M\boxtimes'N$ 的含义同 (260)。比较 (269) 与 (260)（并用基变换同构），即得所要的典范同构 $M\circledast^{c}N\xrightarrow{\sim}M\circledast^{f}N$。

它与结合约束的相容性来自 $X^{3}$ 上的类似图景。写出与交换约束的相容性（用 $\operatorname{Bun}_G$ 与 Hecke）！
% [Tutor Note: 原书 p.209 末句为作者的大写备忘（"WRITE DOWN THE COMAT WITH COM CONSTRAINTS"，即 compatibility with commutativity constraints）；底本如此，保留原样。]

% ===== units/ch05_u03.tex =====
% ch05_u03 —— 原书 p.210–p.219（5.3.18–5.3.31、5.4 节首与 5.4.1–5.4.7）
\SourcePage{210}

\bdp{5.3.18} 对 $M\in\mathcal{P}$，置 $h^{\cdot}(M):=H^{\cdot}_{DR}(\mathcal{GR},M)$。这是 $\mathbb{Z}$-分次向量空间；记 $h^{\varepsilon}(M)$ 为相应的 $\mathbb{Z}/2\mathbb{Z}$-分次向量空间。

考虑投影 $p:\mathcal{GR}_X\to X$。$X$ 上的 $\mathcal{D}$-模 $H^ap_*(M_X)$ 是常值的，即同构于 $\omega_X$ 的若干拷贝之和（回忆：我们处理的是右 $\mathcal{D}$-模）。相应的纤维是 $h^{\cdot}(M)$：对任意 $x\in X$，有 $H^{\cdot}i_x^{!}p_*(M_X)=h^{\cdot}(M)$，其中 $i_x$ 是嵌入 $\{x\}\hookrightarrow X$。

\bdp{5.3.19} \emph{命题.}（[MV]）对 $\mathcal{P}$ 中对象的任意族 $\{M_j\}_{j\in J}$，$X^J$ 上的 $\mathcal{D}$-模 $H^ap_*^{(J)}(\boxtimes M_{jX})$ 是常值的。\hfill$\Box$

对任意 $(x_j)\in X^J$，有
\begin{equation}\tag{270}
H^{\cdot}i^{!}_{(x_j)}p_*^{(J)}(\boxtimes M_{jX})=\otimes h^{\cdot}(M_j).
\end{equation}

由 5.3.18，这对 $(x_j)\in U^{(J)}$ 显然；然后用 5.3.19。
% [Tutor Note: 原书 p.210 此处作「then use 5.3.19」，与正在证明的命题同号，底本如此，保留原样。]

\bdp{5.3.20} 对 $(x_j)\in X\subset X^J$，（270）给出典范同构 $h^{\cdot}(\circledast M_j)=\otimes h^{\cdot}(M_j)$，它显然与「约束」（266）相容。于是
\begin{equation}\tag{271}
h^{\cdot}:\mathcal{P}\to\operatorname{Vect}^{\cdot},\qquad h^{\varepsilon}:\mathcal{P}\to\operatorname{Vect}^{\varepsilon}
\end{equation}
是\emph{张量函子}。这里 $\operatorname{Vect}^{\cdot}$ 是带「超」交换约束的 $\mathbb{Z}$-分次向量空间的张量范畴，$\operatorname{Vect}^{\varepsilon}$ 是 $\mathbb{Z}/2\mathbb{Z}$-分次向量空间的类似张量范畴。

\bdp{5.3.21} 可以扭 $\mathcal{P}$ 上的张量结构，以摆脱超向量空间。为此注意：$\mathcal{P}$ 的对象带典范的 $\mathbb{Z}/2\mathbb{Z}$-分次 $\varepsilon$，它由支撑各分支的奇偶性给出（见 4.5.10）。该分次与 $\boxtimes$ 相容。

记 $\mathcal{P}^{\natural}$ 为 $\mathcal{P}^{\varepsilon}:=\mathcal{P}\otimes\operatorname{Vect}^{\varepsilon}$ 中偶对象的全子范畴（按 $\mathbb{Z}/2\mathbb{Z}$-分次的张量积取）。它是 $\mathcal{P}^{\varepsilon}$ 的张量子范畴。「遗忘分次」函子 $o_{\varepsilon}:\operatorname{Vect}^{\varepsilon}\to\operatorname{Vect}$ 给出等价 $\mathcal{P}^{\natural}\xrightarrow{\sim}\mathcal{P}$。这是幺半范畴的等价（即它与张量积、结合约束相容）；
\SourcePage{211}
$\mathcal{P}$ 与 $\mathcal{P}^{\natural}$ 的交换约束 $A\otimes B\xrightarrow{\sim}B\otimes A$ 相差 $(-1)^{p(A)p(B)}$。
% [Tutor Note: 原书 p.211 起数处出现作者彼此提醒用的大写插语（WHAT ABOUT GINZBURG ??、WHAT ABOUT THE SIGN??? 等），译文按其语气译出，底本如此。]

由 5.3.3(ii)，函子 $h^{\varepsilon}$ 与 $\mathbb{Z}/2\mathbb{Z}$-分次相容。因此它定义张量函子
\begin{equation}\tag{272}
h:\mathcal{P}^{\natural}\to\operatorname{Vect}.
\end{equation}
注意 $h$ 带典范的 $\mathbb{Z}$-分次，因记号滥用，我们仍把它记作 $h^{\cdot}$。于是 $h^{\cdot}$ 是 $\mathcal{P}^{\natural}$ 上的张量函子，取值于带普通（非超）交换约束的分次向量空间张量范畴。

\bdp{5.3.22} 据 [MV]（GINZBURG 怎么办？？），张量范畴 $\mathcal{P}^{\natural}$ 是刚性的，即每个对象都有 [Del91] §2.1.2 意义下的对偶（对偶对象在 [MV] 中有显式构造）。张量函子（272）是 $\mathbb{C}$-线性的且正合\footnote{正合性显然，因为 $\mathcal{P}^{\natural}$ 半单。Mirković 与 Vilonen [MV] 必须证明正合性，因为他们希望其证明对任意交换环上的反常层都有效。}，故它是 [Del91] 意义下的纤维函子。于是由一般的 Tannakian 形式化，（272）诱导张量范畴 $\mathcal{P}^{\natural}$ 与 $\operatorname{Rep}(\operatorname{Aut}^{\otimes}h)$ 之间的等价；这里 $\operatorname{Aut}^{\otimes}h$ 记 $h$ 的张量自同构群概形，$\operatorname{Rep}$ 指有限维表示范畴。据 [MV]，有同构 $\varkappa:{}^{L}G\xrightarrow{\sim}\operatorname{Aut}^{\otimes}h$，故可把上述等价改写为
\begin{equation}\tag{273}
h:\mathcal{P}^{\natural}\xrightarrow{\sim}\operatorname{Rep}{}^{L}G.
\end{equation}
这里 ${}^{L}G$ 是朗兰兹对偶群：它是半单群，带固定的嘉当环面 ${}^{L}H\subset{}^{L}G$、相应的根数据与 $G$ 的根数据之对偶的等同，以及对每个单负根 $\alpha$ 固定的一组非零向量 $y_{\alpha}\in({}^{L}\mathfrak{g})^{\alpha}$。

\SourcePage{212}

\bdp{5.3.23} 我们将通过列举 ${}^{L}G$ 在 $h$ 上作用的若干性质，来定义一个\emph{典范同构}
\begin{equation}\tag{274}
\varkappa:{}^{L}G\xrightarrow{\sim}\operatorname{Aut}^{\otimes}h,
\end{equation}
这些性质唯一确定 $\varkappa$。

(i) 记
\begin{equation}\tag{275}
t\mapsto t^{2\rho}
\end{equation}
为相应于 $G$ 的权 $2\rho$ 的态射 $\mathbb{G}_m\to{}^{L}H$。则 $t^{2\rho}$ 以乘以 $t^{-a}$ 的方式作用在 $h^a$ 上（故单参数子群（275）的作用对应于 $h$ 的分次 $h^{\cdot}$）。由 (i) 推出，${}^{L}H$ 在 $h$ 上的作用保持 $h$ 的分次。

(ii) 对任意 $\chi\in P_+({}^{L}G)$，群 ${}^{L}H$ 通过特征标 $\chi$ 作用在 $h^{\min}(I_{\chi})=h^{-\dim\operatorname{Orb}_{\chi}}(I_{\chi})$ 上。这意味着不可约 ${}^{L}G$-模 $h(I_{\chi})$ 的最高权等于 $\chi$。

\emph{注.} 由于 $\dim\operatorname{Orb}_{\chi}=\langle\chi,2\rho\rangle$，(i) 与 (ii) 之间没有矛盾。

性质 (i) 与 (ii) 可在 [MV] 中找到。它们唯一确定（274）在 ${}^{L}H$ 上的限制。故 (274) 由 (i)、(ii) 唯一确定到 ${}^{L}H$-共轭。我们用下述性质规范化（274）。

(iii) 设 $c\in(\operatorname{Sym}^2\mathfrak{g}^*)^G$ 是 $\mathfrak{g}$ 上的不变双线性型（约化情形下是 $[\mathfrak{g},\mathfrak{g}]$ 上的呢？？？）。置
\begin{equation}\tag{276}
f_c:=\frac{???}{2}\sum_{\alpha}c(\alpha,\alpha)y_{\alpha}\in{}^{L}\mathfrak{g}
\end{equation}
% [Tutor Note: 原书 p.212 公式（276）的分子印作「???」，系底本缺失的记号，底本如此，保留原样。]
（表达式 $c(\alpha,\alpha)$ 有意义，因为 $\alpha\in({}^{L}\mathfrak{h})^*=\mathfrak{h}\subset\mathfrak{g}$。）则李代数元素 $f_c$ 以乘以 $\nu(c)$ 的方式作用在 $h(M)=H^{\cdot}_{DR}(\mathcal{GR},M)$ 上，其中 $M\in\mathcal{P}^{\natural}=\mathcal{P}$，而
\begin{equation}\tag{277}
\nu:(\operatorname{Sym}^2\mathfrak{g}^*)^G\to H^2_{DR}(\mathcal{GR})
\end{equation}

\SourcePage{213}

式（277）是标准态射，其定义将在 5.3.24 中回顾。

\emph{注.} (iii) 由 V. Ginzburg [Gi95] 以略有不同的形式给出。事实上，他用类似方式描述了 $f_c$ 在 ${}^{L}\mathfrak{g}$ 中的整个中心化子在 $h$ 上的作用。

\bdp{5.3.24} 本小节（读者可跳过）定义典范态射（277）。我们采用如下的 ad hoc 定义：对任意归纳概形 $Z$，置 $H^a_{DR}(Z):=\varprojlim H^a(Y,\Omega_Y)$，其中 $Y$ 跑遍 $Z$ 的所有闭子概形，$\Omega_Y$ 是 $Y$ 的 de Rham 复形（最朴素意义下）。为定义 $\nu$，姑且设（简简单单两次？？糟糕的文风）$G$ 半单\footnote{任意 $G$ 的 $\nu$ 之定义留作习题。}。则投影 $G(K)\to\mathcal{GR}$ 诱导同构 $H^2_{DR}(\mathcal{GR})\xrightarrow{\sim}H^2_{DR}(G(K))$（事实上，该投影是 $G(\mathcal{O})$-torsor，$G(\mathcal{O})$ 连通，且 $H^1_{DR}(G(\mathcal{O}))=H^2_{DR}(G(\mathcal{O}))=0$）。现在 $c$ 定义了 $\mathfrak{g}\otimes K$ 上的 Kac–Moody 余圈 $u,v\mapsto\operatorname{Res}_{t=0}c(du,v)$。设 $\omega_c$ 是 $G(K)$ 上相应的右不变闭 2-形式。其类经上述同构之逆映射所得的像为 $\nu(c)\in H^2_{DR}(\mathcal{GR})$。符号怎么办？？？

\emph{注.} 在 5.3.23(iii) 中，我们用了 $H^{\cdot}_{DR}(\mathcal{GR})$ 在 $H^{\cdot}_{DR}(\mathcal{GR},M)$ 上的作用，其中 $M$ 是 $\mathcal{GR}$ 上的 $\mathcal{D}$-模。该作用定义如下。考虑 $\Omega^{\cdot}$-复形 $\Omega M$（见 7.11.13）。则 $H^{\cdot}_{DR}(\mathcal{GR},M)=\varinjlim H^{\cdot}(Y,\Omega M_{(Y)})$，其中 $Y$ 跑遍 $\mathcal{GR}$ 的所有子概形。现在 $\Omega M_{(Y)}$ 是 $Y$ 上的 $\Omega$-复形，故 $H^{\cdot}(Y,\Omega_Y)$ 作用在 $H^{\cdot}(Y,\Omega M_{(Y)})$ 上。因此 $H^{\cdot}_{DR}(\mathcal{GR})$ 作用在 $H^{\cdot}_{DR}(\mathcal{GR},M)$ 上。

\bdp{5.3.25} 5.3.23 中给出的典范同构（274）的简要刻画已足够我们使用。想更好理解（274）的读者可读 ???-??? 与 [MV]。
% [Tutor Note: 原书 p.213 此处作「the reader may read ???-???」，交叉引用印作问号，底本如此，保留原样。]

\bdp{5.3.26}

\SourcePage{214}

\emph{注.} 回忆（见 4.5.9）$\mathcal{GR}$ 的连通分支由 $Z({}^{L}G)^{\vee}$ 的元素标号；这里 $Z({}^{L}G)^{\vee}$ 是中心 $Z({}^{L}G)\subset{}^{L}G$ 的特征标群。$\mathcal{GR}$ 中相应于 $\zeta\in Z({}^{L}G)^{\vee}$ 的连通分支记作 $\mathcal{GR}_{\zeta}$。支撑分解 $D(\mathcal{GR})=\prod D(\mathcal{GR}_{\zeta})$、$\mathcal{P}=\oplus\mathcal{P}_{\zeta}$ 定义了 $h$ 上的 $Z({}^{L}G)^{\vee}$-分次，即 $Z({}^{L}G)$-作用。该作用与 ${}^{L}G$-作用所诱导者一致。

本节余下部分解释上述构造如何与过到 ${}^{L}G$ 的 Levi 子群相容。当该子群为 ${}^{L}H\subset{}^{L}G$ 时，这归结为 Mirković–Vilonen 给出的 ${}^{L}H$ 在纤维函子 $h$ 上作用的显式描述。

\bdp{5.3.27} 设 $P\subset G$ 是抛物子群，$N_P\subset P$ 是其幂幺根，$F:=P/N_P$ 是 Levi 群。$F$ 与 $G$ 的嘉当环面以显然方式等同，$F$ 的根数据是 $G$ 的根数据的子集。于是对标准环面 ${}^{L}H\subset{}^{L}F\subset{}^{L}G$，${}^{L}F$ 是 ${}^{L}G$ 的 Levi 子群。从而 $Z({}^{L}G)\subset Z({}^{L}F)$。

我们将定义典范张量函子
\begin{equation}\tag{278}
r_P^{\natural}:\mathcal{P}_G^{\natural}\to\mathcal{P}_F^{\natural},
\end{equation}
它经由等价 $h_G$、$h_F$ 对应于显然的限制函子 $r^{GF}:\operatorname{Rep}{}^{L}G\to\operatorname{Rep}{}^{L}F$。

\bdp{5.3.28} 图表 $G\leftarrow P\rightarrow F$ 给出相应仿射 Grassmann 流形的态射
\begin{equation}\tag{279}
\mathcal{GR}^G\xleftarrow{i}\mathcal{GR}^P\xrightarrow{\pi}\mathcal{GR}^F.
\end{equation}
这里 $\pi$ 是形式光滑的归纳仿射满投影。其纤维是 $N_P(K)$-轨道。故 $\pi$ 给出 $\mathcal{GR}^P$ 与 $\mathcal{GR}^F$ 的连通分支集之间的双射。对任意 $\zeta\in Z({}^{L}F)^{\vee}$，记 $\mathcal{GR}^P_{\zeta}$ 为相应的分支。则 $i$ 的限制 $i_{\zeta}:\mathcal{GR}^P_{\zeta}\hookrightarrow\mathcal{GR}^G$ 是局部闭嵌入；其像落在 $\mathcal{GR}^G_{\bar{\zeta}}$ 中，其中 $\bar{\zeta}:=\zeta|_{Z({}^{L}G)}$。归纳概形 $\mathcal{GR}^P_{\zeta}$ 构成 $\mathcal{GR}^G$ 的分层（即对任意闭
\SourcePage{215}
子概形 $Y\subset\mathcal{GR}^G$，交 $Y_{\zeta}:=Y\cap\mathcal{GR}^P_{\zeta}$ 构成 $Y$ 的分层）。

置 $\rho_{GF}:=\rho_G-\rho_F\in\mathfrak{h}^*$。由于 $2\rho_{GF}$ 是 $F$ 的特征标（伴随作用在 $\mathfrak{n}_P$ 上的行列式），可把它看作 $Z({}^{L}F)\subset{}^{L}H$ 的单参数子群。于是对如上每个 $\zeta$，有一个整数 $\langle\zeta,2\rho_{GF}\rangle$。记 $\mathcal{GR}^F_n$ 为满足 $\langle\zeta,2\rho_{GF}\rangle=n$ 的分支 $\mathcal{GR}^F_{\zeta}$ 之并。有相应的分解 $D(\mathcal{GR}^F)=\prod D(\mathcal{GR}^F_n)$、$\mathcal{P}^F=\oplus\mathcal{P}^F_n$。置 $\mathcal{P}^{F'}=\oplus\mathcal{P}^F_n[-n]\subset D(\mathcal{GR}^F)$。如同 5.3.18，对 $M\in\mathcal{P}^{F'}$ 置 $h_F(M)=H^{\cdot}(\mathcal{GR}^F,M)\in\operatorname{Vect}^{\cdot}$。

\bdp{5.3.29} \emph{命题.}
\begin{enumerate}[label=(\roman*),leftmargin=2.2em]
\item 函子 $r_D^{GF}:=\pi_*i^{!}:D(\mathcal{GR}^G)\to D(\mathcal{GR}^F)$ 把 $\mathcal{P}^G$ 送入 $\mathcal{P}^{F'}$，故有
\begin{equation}\tag{280}
r_P^{GF}:\mathcal{P}^G\to\mathcal{P}^{F'}.
\end{equation}
\item 有函子的典范等同
\begin{equation}\tag{281}
h_G^{\cdot}=h_F^{\cdot}r_P^{GF}:\mathcal{P}^G\to\operatorname{Vect}^{\cdot}.
\end{equation}
\end{enumerate}

\emph{证明.} 先设 $P=B$ 是 Borel 子群。则 $F=H$，且 $\mathcal{GR}^H_{red}=({}^{L}H)^{\vee}$，故 $\mathcal{GR}^H$ 上的 $\mathcal{D}$-模等同于 $({}^{L}H)^{\vee}$-分次向量空间，即 ${}^{L}H$-模。分层 $\mathcal{GR}^B_{\zeta}$ 正是 $\mathcal{GR}^G$ 上的 $N_B(K)$-轨道。因此 5.3.29 就是 [MV] 的关键定理。

回忆等同（281）的构造如下（见 [MV]）。设 $\overline{\mathcal{GR}}^B_n\subset\mathcal{GR}^G$ 为 $\mathcal{GR}^B_n:=\pi^{-1}(\mathcal{GR}^H_n)$ 在 $\mathcal{GR}^G$ 中的闭包。则各 $\overline{\mathcal{GR}}^B_n$ 给出 $\mathcal{GR}^G$ 上的递减滤链。对任意 $M\in\mathcal{P}^G$，显然的态射
\[
h^n_Hr_P^{GH}(M)=H^n(\mathcal{GR}^B_n,i^{!}M)\leftarrow H^n_{\overline{\mathcal{GR}}^B_n}(\mathcal{GR}^G,M)\rightarrow H^n(\mathcal{GR}^G,M)=h^n_G(M)
\]
都是同构。它们的复合是（281）。

现在设 $P$ 是任意抛物子群。取 Borel 子群 $B\subset P$，于是 $B_F:=B/N_P\cap B$ 是 $F$ 的 Borel 子群。考虑函子 $r_D^{GH}:D(\mathcal{GR}^G)\to D(\mathcal{GR}^H)$ 与 $r_D^{FH}:D(\mathcal{GR}^F)\to D(\mathcal{GR}^H)$。由基变换，有函子的典范等同 $r_D^{GH}=r_D^{FH}r_D^{GF}$。设
\SourcePage{216}
$\mathcal{P}^{H'}\subset D(\mathcal{GR}^H)$ 是由 $B\subset G$ 所定义的范畴，故已知 $r_D^{GH}(\mathcal{P}^G)\subset\mathcal{P}^{H'}$，且（因 $\rho_{GF}=\rho_G-\rho_F$）有 $r_D^{FH}(\mathcal{P}^{F'})\subset\mathcal{P}^{H'}$。

函子 $r_P^{GH}:\mathcal{P}^{F'}\to\mathcal{P}^{H'}$ 忠实（因为差一个平移，它与 $h_F^{\cdot}$ 一致）。于是使所有 $H^iT$ 都属于 $\mathcal{P}^F$ 的对象 $T\in D(\mathcal{GR}^F)$ 属于 $\mathcal{P}^{F'}$，当且仅当 $r_D^{FH}(T)\subset\mathcal{P}^{H'}$。
% [Tutor Note: 原书 p.216 此处作「since up to shift if coincides with h^·_F」，其中「if」疑为「it」之误，底本如此，保留原样。]

把该注记用于 $T=r_D^{GF}(M)$（$M\in\mathcal{P}^G$），可见 $r_D^{GF}(M)\in\mathcal{P}^{F'}$，这正是 5.3.29(i)。我们还知道 $h_G^{\cdot}(M)=h_H^{\cdot}(r_P^{GH}(M))=h_H^{\cdot}(r_P^{FH}(M))=h_F^{\cdot}r^{GF}(M)$，这就是等同 5.3.29(ii)。至于它不依赖辅助选取的 Borel 子群 $B\subset P$，留作习题。\hfill$\Box$

\bdp{5.3.30} 范畴 $\mathcal{P}^{F'}$ 带典范张量结构（由用于 $\mathcal{P}^F$ 的同样构造定义）。函子 $r_P^{GF}:\mathcal{P}^G\to\mathcal{P}^{F'}$ 以典范方式为张量函子。事实上，（279）是手征半群的态射，故可考虑相应的函子 $r_D^{GF}:=\pi_*i^{!}:D(\mathcal{GR}^G_{X^J})\to D(\mathcal{GR}^F_{X^J})$。留作习题（提示：用 5.3.19）：对 $M_j\in\mathcal{P}^G$，该函子把 $\boxtimes M_j$ 送到 $\boxtimes r_D^{GF}(M_j)$（记号见 5.3.14）。由于它（由基变换）还与函子 $\widetilde{\Delta}^{(J)!}$ 交换，我们得到所需的张量积相容性。如同 5.3.19 可见，（281）是张量函子的同构。

最后，如同 5.3.21，把张量范畴 $\mathcal{P}^G$ 换成 $\mathcal{P}^{G\natural}$。由于 $\rho_{GF}=\rho_G-\rho_F$，可见 $r_P^{GF}$ 给出张量函子 $r^{GF}:\mathcal{P}^{G\natural}\to\mathcal{P}^{F\natural}$，它与纤维函子 $h_G$、$h_F$ 相容。它定义了态射 $r:\operatorname{Aut}^{\otimes}h_F\to\operatorname{Aut}^{\otimes}h_G$。

\bdp{5.3.31} \emph{引理.} 态射 $\varkappa_G^{-1}r\varkappa_F:{}^{L}F\to{}^{L}G$ 与 5.3.27 中的典范嵌入一致。\hfill$\Box$

\bdsec{5.4. 主定理 II：从局部到整体}

本节给出 5.2 中主定理的精确版本，并证明局部主定理蕴含整体主定理。我们本质地用到第 7 章的「Hecke 模式」。要理解发生了什么，必须（且几乎充分）阅读 7.1.1 与 7.9.1。

\SourcePage{217}

\bdp{5.4.1} 我们从 Hecke 特征 $\mathcal{D}$-模的定义开始。考虑带 $\operatorname{Aut}\mathcal{O}$ 作用的对 $(G(K),G(\mathcal{O}))$。设 $\mathcal{H}$ 是相应的 $(\operatorname{Der}\mathcal{O},\operatorname{Aut}^0\mathcal{O})$-等变 Hecke 范畴，定义见 7.9.2\footnote{我们的 $(G(K),G(\mathcal{O}))$、$(\operatorname{Der}\mathcal{O},\operatorname{Aut}^0\mathcal{O})$ 分别是 7.9.2 的 $(G,K)$、$(\mathfrak{l},P)$。}。由于 $\mathcal{P}$ 的每个对象都以典范方式为 $\operatorname{Aut}\mathcal{O}$-等变 $\mathcal{D}$-模\footnote{据 5.3.4，$\mathcal{P}$ 的每个对象带唯一的强 $\operatorname{Aut}^0\mathcal{O}$-作用，它与强 $\operatorname{Aut}\mathcal{O}$-作用相同。}，$\mathcal{P}$ 是 $\mathcal{H}$ 的全子范畴。由定义可知，嵌入 $\mathcal{P}\to\mathcal{H}$ 是幺半函子。

考虑 $X$ 上典范的 $\operatorname{Aut}\mathcal{O}$-结构 $X^{\wedge}$（见 2.6.5）与 2.8.3 中定义的 $X^{\wedge}$ 上的概形 $M^{\wedge}$；它带 $\operatorname{Aut}\mathcal{O}\times G(K)$ 的典范作用（见 2.8.3–2.8.4）。商叠 $(\operatorname{Aut}^0\mathcal{O}\times G(\mathcal{O}))\backslash M^{\wedge}$ 等于 $\operatorname{Bun}_G\times X$。于是我们来到 7.9.1、7.9.4\footnote{我们的 $X^{\wedge}$ 与 $M^{\wedge}$ 是 7.9.4 的 $X^{\wedge}$ 与 $Y^{\wedge}$。} 的设定。故 $\mathcal{H}$ 作用在 $D(\operatorname{Bun}_G\times X)$ 上。因此 $D(\operatorname{Bun}_G\times X)$ 是 $\mathcal{P}$-模。借助 Satake 等价（273）把幺半范畴 $\mathcal{P}$\footnote{本节（5.4.6 的注除外）只用 $\mathcal{P}$ 的幺半结构（交换约束不起作用）。故可把 $\mathcal{P}$ 与 $\mathcal{P}^{\natural}$ 等同。} 与 $\operatorname{Rep}{}^{L}G$ 等同，便得到 $\operatorname{Rep}{}^{L}G$ 在 $D(\operatorname{Bun}_G\times X)$ 上的典范作用，称为 Hecke 作用。记作 $\circledast$。

注意 $D(\operatorname{Bun}_G\times X)$ 还带张量范畴 $\operatorname{Vect}^{\nabla}(X)$ 的显然作用；$\operatorname{Vect}^{\nabla}(X)$ 由 $X$ 上带联络的向量丛构成（事实上，由 $X$ 上无挠左 $\mathcal{D}$-模构成的更大张量范畴亦可）。我们把它记作 $\otimes$。它与 Hecke 作用交换，故 $D(\operatorname{Bun}_G\times X)$ 是 $(\operatorname{Rep}{}^{L}G,\operatorname{Vect}^{\nabla}(X))$-双模。

设 $\mathfrak{F}$ 是 $X$ 上带联络的 ${}^{L}G$-丛。它给出张量函子 $\operatorname{Rep}{}^{L}G\to\operatorname{Vect}^{\nabla}(X)$，$V\mapsto V_{\mathfrak{F}}$，从而给出 $\operatorname{Rep}{}^{L}G$ 在 $D(\operatorname{Bun}_G\times X)$ 上相应的作用。

\bdp{5.4.2} 设 $M$ 是 $\operatorname{Bun}_G$ 上的 $\mathcal{D}$-模。设 $M_{(X)}\in\mathcal{M}(\operatorname{Bun}_G\times X)$ 是 $M$ 的拉回。假设对任意 $V\in\operatorname{Rep}{}^{L}G$，给定自然的
\SourcePage{218}
同构 $\alpha_V:V\circledast M_{(X)}\xrightarrow{\sim}M_{(X)}\otimes V_{\mathfrak{F}}$（特别地，$V\circledast M_{(X)}$ 是 $\mathcal{D}$-模，而不仅是导出范畴的对象）。若对任意 $V_1,V_2\in\operatorname{Rep}{}^{L}G$ 有 $\alpha_{V_1\otimes V_2}=\alpha_{V_1}\circ(V_1\circledast\alpha_{V_2})$，则称这些 $\alpha_V$ 在 $M$ 上定义了\emph{Hecke $\mathcal{F}$-特征模结构}。称这样的 $(M,\alpha_V)$（或简记 $M$）为\emph{Hecke $\mathfrak{F}$-特征模}。
% [Tutor Note: 原书 p.218 前两处 oper 记号印作花体 ℱ，同页后文与 5.4.3 中均印作 𝔉，底本混用，保留原样。]

\emph{注.} 对 $X$ 上任意 ${}^{L}G$-局部系统 $\mathcal{F}$，我们希望定义 Hecke $\mathfrak{F}$-特征模的三角范畴\footnote{当然，在 Hecke 特征模的上述定义中，$M$ 可取 $D(\operatorname{Bun}_G)$ 的任意对象，而不只是 $\mathcal{D}$-模。但在此一般性下该定义看来并不合理（这类对象不构成三角范畴）。}。

下述定理是定理 5.2.6 的精确版本。

\bdp{5.4.3} \emph{定理.} 对任意 ${}^{L}G$-oper $\mathfrak{F}$，5.1.1 中定义的 $\mathcal{D}$-模 $M_{\mathfrak{F}}$ 有自然的 Hecke $\mathfrak{F}$-特征模结构。

留作习题：函子 $T^i_X$ 与 $H^iV^X\circledast$ 一致（记号见 5.2.4、5.2.5）。于是定理 5.4.3 蕴含 5.2.6。

\bdp{5.4.4} 我们需要 5.4.1–5.4.3 的「带参数」版本。设 $A$ 是交换环。记 $\mathcal{M}(\operatorname{Bun}_G\times X,A)$ 为 $\mathcal{M}(\operatorname{Bun}_G\times X)$ 中 $A$-模构成的范畴（即带 $A$-作用的 $\mathcal{D}$-模）。它有导出版本 $D(\operatorname{Bun}_G\times X,A)$，后者是以 $\mathcal{M}(\operatorname{Bun}_G\times X,A)$ 为心的 t-范畴（见 7.3.13）。如同 5.4.1，范畴 $D(\operatorname{Bun}_G\times X,A)$ 带 $\operatorname{Rep}{}^{L}G$ 的 Hecke 作用。

我们还有 $A\otimes\mathcal{O}_X$-平坦的 $A\otimes\mathcal{D}_X$-模所成张量范畴在 $D(\operatorname{Bun}_G\times X,A)$ 上的显然作用，它与 Hecke 作用交换。因此 $X$ 上带联络的 ${}^{L}G$-丛的任意平坦 $A$-族 $\mathcal{F}_A$ 给出 $\operatorname{Rep}{}^{L}G$ 在 $D(\operatorname{Bun}_G\times X,A)$ 上的作用。

现在对 $M\in\mathcal{M}(\operatorname{Bun}_G,A)$，如同 5.4.2 定义 $M$ 上 Hecke $\mathfrak{F}_A$-特征模结构的概念。下述定理是 5.2.9 的精确版本；由 5.1.2(i)，它蕴含 5.4.3。

\SourcePage{219}

\bdp{5.4.5} \emph{定理.} 5.1.1 中定义的 $\mathcal{D}$-模 $M_{\mathcal{L}}\in\mathcal{M}(\operatorname{Bun}_G,\mathcal{A}_{{}^{L}\mathfrak{g}}(X))$ 有典范的 Hecke $\mathfrak{F}_{\mathcal{L}}$-特征模结构。

\bdp{5.4.6} \emph{注.} (i) 有时（当你想用交换约束时，例如见下一个注或下一节）在超 $\mathcal{D}$-模的框架下处理上述概念是方便的。注意 $\operatorname{Bun}_G$ 上任意 $\mathcal{D}$-模 $M$ 都带典范的 $\mathbb{Z}/2\mathbb{Z}$-分次：$M$ 为偶或奇，取决于 $M$ 支撑在 $\operatorname{Bun}_G$ 的偶分支还是奇分支。记该超 $\mathcal{D}$-模为 $M^{\natural}$。于是 $\natural$ 把 $\mathcal{M}(\operatorname{Bun}_G)$ 与 $\mathcal{M}(\operatorname{Bun}_G)^{\varepsilon}:=\mathcal{M}(\operatorname{Bun}_G)\otimes\operatorname{Vect}^{\varepsilon}$ 的全子范畴 $\mathcal{M}(\operatorname{Bun}_G)^{\natural}$ 等同。对 $D(\operatorname{Bun}_G)$ 与 $D(\operatorname{Bun}_G\times X)$ 同样适用。

$\mathcal{P}$ 在 $D(\operatorname{Bun}_G\times X)$ 上的作用给出 $\mathcal{P}^{\varepsilon}$ 在 $D(\operatorname{Bun}_G\times X)^{\varepsilon}$ 上的作用。$\mathcal{P}^{\natural}\subset\mathcal{P}^{\varepsilon}$ 的作用保持 $D(\operatorname{Bun}_G\times X)^{\natural}$，也保持 $\operatorname{Vect}^{\nabla}(X)$-作用。现在完全如同 5.4.2 定义 $\mathcal{M}(\operatorname{Bun}_G)^{\natural}$ 的 Hecke $\mathfrak{F}$-特征对象的概念。该定义没有带来新东西：$\mathcal{D}$-模 $M$ 是 Hecke $\mathfrak{F}$-特征模，当且仅当 $M^{\natural}$ 是。

(ii) 上述 $\mathcal{M}(\operatorname{Bun}_G)$ 中 $M$ 上的 $\mathcal{F}$-特征模结构定义，用了 $\mathcal{P}$ 上张量结构的卷积构造。可以改用 $\circledast$ 的融合构造重写如下。DOPISAT'!!!
% [Tutor Note: 原书 p.219 此处为俄文注记「DOPISAT'!!!」，意为「待补写」，底本如此，保留原样。]

\bdp{5.4.7} 转向 5.2 中的主要局部定理。我们处于 5.2.12 的设定，故固定 $\mathcal{L}\in Z\operatorname{tors}_{\theta}(\mathcal{O})$，它定义了 $G(K)$ 的中心扩张 $\widetilde{G(K)}=\widetilde{G(K)}_{\mathcal{L}}$，该扩张在群子概形 $G(\mathcal{O})$ 上分裂（见 4.4.9）。我们有相应的扭 Harish-Chandra 模范畴 $\mathcal{M}(\mathfrak{g}\otimes K,G(\mathcal{O}))'$ 与 Harish-Chandra 复形的导出范畴 $D(\mathfrak{g}\otimes K,G(\mathcal{O}))'$（见 7.8.1 与 7.14.1\footnote{故 $1\in\mathbb{C}\subset\mathfrak{g}\otimes K$ 在这些范畴的对象上通过恒等作用。}）。据 7.8.2、7.14.1，$D(\mathfrak{g}\otimes K,G(\mathcal{O}))'$ 带对 $(G(K),G(\mathcal{O}))$ 的 Hecke 幺半范畴 $\mathcal{H}$ 的典范作用 $\circledast$。由于 $\mathcal{P}$ 是 $\mathcal{H}$（的心）的幺半子范畴，我们的 $D(\mathfrak{g}\otimes K,G(\mathcal{O}))'$ 是 $\mathcal{P}$-模。
% 原书 p.219 末句完结，但 5.4.7 的论述续于 p.220（下一单元）。

% ===== units/ch05_u04.tex =====
% ch05_u04 —— 原书 p.220--p.229（5.4.8 -- 5.6.11）
% 本单元从 5.4.8 起，末尾停在 5.6.11 的中途；尾部接 p.230。

\SourcePage{220}

设 $\mathit{Vac}'\in\mathcal{M}(\mathfrak{g}\otimes K,G(O))'$ 为扭真空模。

\bdp{5.4.8} \emph{定理.} 对任意对象 $P\in\mathcal{P}$，对象
$P\circledast\mathit{Vac}'\in D(\mathfrak{g}\otimes K,G(O))'$ 同构于若干份
$\mathit{Vac}'$ 的直和。\footnote{特别地，它是单个 Harish-Chandra 模，而非这样的复形。}

这一定理等价于 5.2.14。事实上，由 7.8.5 的（335）与 7.14.1，有
$(\widetilde{\mathfrak{g}\otimes K},G(O))$-模的典范等同
\begin{equation}\tag{282}
H^{i}(P\circledast\mathit{Vac}')=H^{i}(\mathcal{GR},P\lambda_{\mathcal{L}}^{-1}).
\end{equation}

这里 $P\lambda_{\mathcal{L}}^{-1}:=P\otimes\lambda_{\mathcal{L}}^{-1}$。有兴趣的读者可以
直接去看这一定理的证明，它可在 ??? 中找到。
% [Tutor Note: 原书 p.220 此处交叉引用印作「???」，底本未补，译文照录。]

\bdp{5.4.9} 我们需要把 $\operatorname{Aut}O$ 的对称性纳入上述设定。回忆（见 4.6.6）
$\operatorname{Aut}O$ 在 $G(K)$ 上的作用提升为 $\operatorname{Aut}_ZO$ 在
$\widetilde{G(K)}$ 上的作用，后者保持 $G(O)$。所以我们处在 7.9.5 的设定中。
\footnote{我们的 $(\operatorname{Der}O,\operatorname{Aut}^{0}_{Z}O)$ 与
$(\widetilde{G(K)},G(O))$ 就是 7.9.5 的 $(\mathfrak{l},P)$ 与 $(G',K)$。}

设 $D_{HC}$ 为 7.9.5 中定义的 Harish-Chandra 复形的导出范畴。这是一个
t-范畴，其心 $\mathcal{M}_{HC}$ 等于对
$(\operatorname{Der}O\ltimes\widetilde{\mathfrak{g}\otimes K},
\operatorname{Aut}^{0}_{Z}O\ltimes G(O))$ 的 Harish-Chandra 模的范畴
（我们假设中心 $\mathfrak{C}\subset\tilde{\mathfrak{g}}\otimes K$ 按标准方式作用）。
% [Tutor Note: 原书 p.220 此处「中心」印作 fraktur C（$\mathfrak{C}$），按上下文
% 应为 $\mathfrak{Z}$；底本如此，译文照录。]

对 $(G(K),G(O))$ 的 $(\operatorname{Der}O,\operatorname{Aut}^{0}_{Z}O)$-等变
Hecke 范畴（见 7.9.2）以全幺半子范畴的形式包含相应的
$(\operatorname{Der}O,\operatorname{Aut}^{0}O)$-等变范畴 $\mathcal{H}$ 与
$\mathcal{H}^{c}$。于是由 7.9.5，$D_{HC}$ 是 $\mathcal{H}$-模，从而是
$\mathcal{P}$-模。

我们需要对设定稍作改变。按惯例，设 $\mathfrak{Z}$ 为 $\mathfrak{g}\otimes K$ 的
完备化扭泛包络代数的中心，$\mathfrak{z}$ 为扭真空模 $\mathit{Vac}'$ 的自同态环；
有显然的代数态射 $e:\mathfrak{Z}\to\mathfrak{z}$。设 $D_{HC\mathfrak{z}}$ 为 7.9.8 中
定义的相应 Harish-Chandra 复形的导出范畴（另见
\SourcePage{221}
7.9.7(iii)）\footnote{我们的 $D_{HC\mathfrak{z}}$ 就是 7.9.8 的 $D^{e}_{HC\,A}$。
在 7.9.8 中，$\mathfrak{Z}$ 表示中心里 $G(K)$-不变的元素之集，但由 3.7.7(ii)，
中心的所有元素都是 $G(K)$-不变的。}。这是一个 t-范畴，其心
$\mathcal{M}_{HC\mathfrak{z}}$ 等于被 $\operatorname{Ker}e$ 消灭的
Harish-Chandra 模的范畴。

设 $\mathcal{H}_{\mathfrak{z}}$ 是 $(G(K),G(O))$ 的
$(\operatorname{Der}O,\operatorname{Aut}^{0}_{Z}O)$-等变 Hecke 范畴的
$\mathfrak{z}$-线性版本，定义见 7.9.7(i)。由 7.9.8，它作用在
$D_{HC\mathfrak{z}}$ 上。由于有显然的幺半函子
$\mathcal{H}\to\mathcal{H}_{\mathfrak{z}}$（见 7.9.7 的注），
$\mathcal{H}_{\mathfrak{z}}$ 包含 $\mathcal{P}$，故 $D_{HC\mathfrak{z}}$ 是
$\mathcal{P}$-模。如同 5.4.1，我们借助 Satake 等价把 $\mathcal{P}$ 换成
$\operatorname{Rep}{}^{L}G$，并把 $\operatorname{Rep}{}^{L}G$ 在
$D_{HC\mathfrak{z}}$ 上的相应作用记作 $\circledast$。另一方面，
$\mathcal{H}_{\mathfrak{z}}$ 的中心包含张量范畴
$\mathcal{M}(\operatorname{Aut}_ZO)^{fl}_{\mathfrak{z}}$，它由平坦
$\mathfrak{z}$-模配上 $\operatorname{Aut}_ZO$-作用构成（见 7.9.7(i)）。
$\mathcal{M}(\operatorname{Aut}_ZO)^{fl}_{\mathfrak{z}}$ 在 $D_{HC\mathfrak{z}}$ 上的
相应作用是显然的：对 $W\in\mathcal{M}(\operatorname{Aut}_ZO)^{fl}_{\mathfrak{z}}$、
$V\in D_{HC\mathfrak{z}}$，有
$W\circledast V=W\otimes V:=W\otimes_{\mathfrak{z}}V$。因此
$D_{HC\mathfrak{z}}$ 是
$(\operatorname{Rep}{}^{L}G,\mathcal{M}(\operatorname{Aut}_ZO)^{fl}_{\mathfrak{z}})$-双模。

设 $\mathfrak{F}$ 是 $\operatorname{Spec}\mathfrak{z}$ 上
$\operatorname{Aut}_ZO$-等变的 ${}^{L}G$-torsor。它给出张量函子
$\operatorname{Rep}{}^{L}G\to\mathcal{M}(\operatorname{Aut}_ZO)^{fl}_{\mathfrak{z}}$，
$V\mapsto V_{\mathfrak{F}}$，从而给出 $\operatorname{Rep}{}^{L}G$ 在
$D_{HC\mathfrak{z}}$ 上的相应作用。

\bdp{5.4.10} 我们在当前的 Harish-Chandra 设定中重复 5.4.2 的定义。即，
\emph{Hecke $\mathfrak{F}$-特征模}是一个 Harish-Chandra 模
$M\in\mathcal{M}_{HC\mathfrak{z}}$，连同自然同构
$\alpha_V:V\circledast M\xrightarrow{\sim}M\otimes V_{\mathfrak{F}}$
（$V\in\operatorname{Rep}{}^{L}G$），使得对任意
$V_1,V_2\in\operatorname{Rep}{}^{L}G$ 有
$\alpha_{V_1\otimes V_2}=\alpha_{V_1}\circ(V_1\circledast\alpha_{V_2})$。

现在可以表述 5.2.16 的精确版本。如同 5.2.15，我们的
$\mathcal{L}\in Z\operatorname{tors}_{\theta}(O)$（见 5.4.7）定义了局部
${}^{L}\mathfrak{g}$-oper 的模概形上的 $\operatorname{Aut}_ZO$-等变
${}^{L}G$-torsor\footnote{$\operatorname{Aut}_ZO$ 的作用来自等同
$\operatorname{Aut}_ZO=\operatorname{Aut}(O,\mathcal{L})$；见 4.6.6。}。
借助 Feigin--Frenkel 同构（80）把该概形与 $\operatorname{Spec}\mathfrak{z}$ 等同，
就得到 $\operatorname{Spec}\mathfrak{z}$ 上相应的
$\operatorname{Aut}_ZO$-等变 torsor $\mathfrak{F}_{\mathcal{L}}$。

\SourcePage{222}

从现在起，我们把 $\mathit{Vac}'$ 看作 $\mathcal{M}_{HC\mathfrak{z}}$ 的对象
（相对于固定真空向量的 $\operatorname{Aut}_ZO$-作用）。

\bdp{5.4.11} \emph{定理.} $\mathit{Vac}'$ 具有 Hecke
$\mathfrak{F}_{\mathcal{L}}$-特征模的典范结构。

这一定理蕴含 5.2.16。事实上，同构（282）是
$\operatorname{Aut}_ZO$-等变的，因为 $\operatorname{Aut}_ZO$ 通过结构的传输
作用在（282）的两端。

它将在何处证明？？？
% [Tutor Note: 原书 p.222 此处作「Where will it be proved???」，三个问号系底本
% 如此（作者未注明证明所在），译文照录。]

现在可以转向本节的主要结果。

\bdp{5.4.12} \emph{定理.} 定理 5.4.11 蕴含 5.4.5。

\emph{证明.} 我们将说明，一个适当的「局部化函子」${}^{L}\Delta$ 把局部图景
变为整体图景\footnote{读者可自行判断，这种疯狂之中是否有章法。}。

我们需要对 5.4.1 的设定稍作修改，以便使用 7.9 的「等变 Hecke 模式」。回忆在
整体定理 5.4.5 的表述中，我们固定了
$\mathcal{L}^{\mathrm{glob}}\in Z\operatorname{tors}_{\theta}(X)$（见 5.2.8），而在局部定理
5.4.11 中用的是 $\mathcal{L}^{\mathrm{loc}}\in Z\operatorname{tors}_{\theta}(O)$。考虑
4.4.15 中与 $\mathcal{L}^{\mathrm{glob}}$ 和 $\mathcal{L}^{\mathrm{loc}}$ 相应的概形
$X^{\wedge}_{Z}$ 与 $M^{\wedge}_{Z}$（它们是 5.4.1 所用概形 $X^{\wedge}$ 与
$M^{\wedge}$ 的平展 $Z$-覆叠）。回忆 $\operatorname{Aut}_ZO$ 作用在
$X^{\wedge}_{Z}$ 上，$\operatorname{Aut}_ZO\ltimes G(K)$ 作用在
$M^{\wedge}_{Z}$ 上（见 4.4.15）。有
$\operatorname{Aut}^{0}_{Z}O\backslash X^{\wedge}_{Z}=X$，且商叠
$(\operatorname{Aut}^{0}_{Z}O\ltimes G(O))\backslash M^{\wedge}_{Z}$ 等于
$\operatorname{Bun}_G\times X$。显然，在 5.4.1 中构造
$D(\operatorname{Bun}_G\times X)$ 上的 Hecke 作用时，可以把
$(M^{\wedge},\operatorname{Aut}O\ltimes G(K))$ 换成
$(M^{\wedge}_{Z},\operatorname{Aut}_ZO\ltimes G(K))$。

如同 5.1.1，设 $\lambda_{\mathcal{L}^{\mathrm{glob}}}$ 是 $\operatorname{Bun}_G$ 上对应于
$\mathcal{L}^{\mathrm{glob}}$ 的 Pfaffian 线丛。记
$\hat{\lambda}=\hat{\lambda}_{\mathcal{L}^{\mathrm{glob}}}$ 为它到 $M^{\wedge}_{Z}$ 的拉回。
$\operatorname{Aut}_ZO\ltimes G(K)$ 在 $M^{\wedge}_{Z}$ 上的作用以典范方式提升为
中心扩张 $\operatorname{Aut}_ZO\ltimes\widetilde{G(K)}$ 在 $\hat{\lambda}$ 上的作用
（见 4.4.16）。所以我们处于
\SourcePage{223}
7.9.6 的设定\footnote{很抱歉记号如此不一致：我们的 $M^{\wedge}_{Z}$、$X^{\wedge}$、
$\hat{\lambda}$、$\operatorname{Der}O$、$\operatorname{Aut}^{0}_{Z}O$、
$\widetilde{G(K)}$、$G(O)$ 分别是 7.9.6 的 $Y^{\wedge}$、$X^{\wedge}$、
$\mathcal{L}^{*}$、$\mathfrak{l}$、$P$、$G'$、$K$。}，于是有右 t-正合的
局部化函子
\[
{}^{L}\Delta:\ D_{HC}\to D(\operatorname{Bun}_G\times X)
\]

在 $\mathfrak{z}$-模的设定中也有相应的图景。事实上，按 7.9.7(ii)，考虑
$\mathcal{D}_X$-代数 $\mathfrak{z}_X$\footnote{任意 $\operatorname{Aut}O$-模 $V$
给出 $\mathcal{D}_X$-模 $V_X$，见 2.6.6。}（我们在 2.7 中已经用过它）以及相应的
范畴 $D(\operatorname{Bun}_G\times X,\mathfrak{z}_X)$，后者是
$\operatorname{Bun}_G\times X$ 上带 $\mathfrak{z}_X$-作用的 $\mathcal{D}$-模的
导出范畴（见 7.3.13）。它带有 $\mathcal{H}_{\mathfrak{z}}$ 的典范作用。有典范的
局部化函子
\[
{}^{L}\Delta_{\mathfrak{z}}:\ D_{HC\mathfrak{z}}\to
D(\operatorname{Bun}_G\times X,\mathfrak{z}_X)
\]
它是 $\mathcal{H}_{\mathfrak{z}}$-模的态射。上述 ${}^{L}\Delta$ 彼此相容（它们
与忘掉 $\mathfrak{z}$-作用可交换）。
% [Tutor Note: 原书 p.223 两处印作「Morhism」，系「Morphism」之误，译文按其本意
% 译为「态射」。]

现在我们的定理是下述事实的直接推论：

\begin{enumerate}[label=(\alph*),leftmargin=2.6em]
\item 有自然的等同
\begin{equation}\tag{283}
{}^{L}\Delta(\mathit{Vac}')=\Delta(\mathit{Vac}')
=M_{\mathcal{L}^{\mathrm{glob}}}\otimes\mathcal{O}_X
\end{equation}
使得 $\Delta(\mathit{Vac}')=\Delta_{\mathfrak{z}}(\mathit{Vac}')$ 上的
$\mathfrak{z}_X$-作用，与 $\mathfrak{z}_X$ 经由极大常值商
$\mathfrak{z}(X)\otimes\mathcal{O}_X=A_{{}^{L}\mathfrak{g}}(X)\otimes\mathcal{O}_X$
以及 $M_{\mathcal{L}^{\mathrm{glob}}}$ 上标准的 $A_{{}^{L}\mathfrak{g}}(X)$-模结构在
$M_{\mathcal{L}^{\mathrm{glob}}}\otimes\mathcal{O}_X$ 上的作用一致。证明见 7.14.9
（注意 $\mathfrak{z}_X$ 按结构的传输作用）。
\item 函子 ${}^{L}\Delta_{\mathfrak{z}}$ 是
$(\operatorname{Rep}{}^{L}G,\mathcal{M}(\operatorname{Aut}_ZO)^{fl}_{\mathfrak{z}})$-双模
的态射。事实上，它是 $\mathcal{H}_{\mathfrak{z}}$-模的态射。
\item 对任意 $W\in\mathcal{M}(\operatorname{Aut}_ZO)^{fl}_{\mathfrak{z}}$、
$T\in D(\operatorname{Bun}_G\times X,\mathfrak{z}_X)$，有
$W\circledast T=W_X\otimes_{\mathfrak{z}_X}T$，其中 $W_X$ 是对应于 $W$ 的
$\mathfrak{z}_X$-模。证明见 7.9.7(i)。
\end{enumerate}

\SourcePage{224}

\begin{enumerate}[label=(\alph*),leftmargin=2.6em]
\item[(d)] 对任意 $V\in\operatorname{Rep}{}^{L}G$，有典范等同
\[
(V_{\mathfrak{F}_{\mathcal{L}^{\mathrm{loc}}}})_X\otimes_{\mathfrak{z}_X}
(\mathfrak{z}(X)\otimes\mathcal{O}_X)\xrightarrow{\sim}
V_{\mathfrak{F}_{\mathcal{L}^{\mathrm{glob}}}}
\]
它与 $V$ 的张量积相容（这里 $\mathfrak{F}_{\mathcal{L}^{\mathrm{loc}}}$ 就是 5.4.10 的
$\mathfrak{F}_{\mathcal{L}}$）。\hfill$\Box$
\end{enumerate}

\bdsec{5.5. oper 的诞生}

本节假设定理 5.4.8 成立。我们首先说明，该定理蕴含 $\mathit{Vac}'$ 是 Hecke
$\mathcal{F}$-特征模，其中 $\mathcal{F}$ 是 $\operatorname{Spec}\mathfrak{z}$ 上
某个 $\operatorname{Aut}_ZO$-等变的 ${}^{L}G$-torsor。本节要点在于：
$\mathcal{F}$ 自然地来自局部 oper 的某个 $\operatorname{Aut}_ZO$-等变的
$\mathfrak{z}$-族。稍后我们将看到，
从 $\operatorname{Spec}\mathfrak{z}$ 到局部 oper 模空间的相应映射与
Feigin--Frenkel 同构一致，由此得到局部主定理。

\bdp{5.5.1} 对任意 $V\in\operatorname{Rep}{}^{L}G$，置
\begin{equation}\tag{284}
F_{\mathcal{H}}(V):=\operatorname{Hom}_{\widetilde{\mathfrak{g}\otimes K}}
(\mathit{Vac}',V\circledast\mathit{Vac}')=(V\circledast\mathit{Vac}')^{G(O)}.
\end{equation}

这是 $\operatorname{Aut}_ZO$-等变的 $\mathfrak{z}$-模\footnote{$V\circledast
\mathit{Vac}'$ 上的两个 $\mathfrak{z}$-模结构相同，因为 Hecke 函子是
$\mathfrak{z}$-线性的。}。由 5.4.8，它是自由 $\mathfrak{z}$-模，故
$F_{\mathcal{H}}(V)\in\mathcal{M}(\operatorname{Aut}_ZO)^{fl}_{\mathfrak{z}}$。
有典范同构
\begin{equation}\tag{285}
V\circledast\mathit{Vac}'=\mathit{Vac}'\otimes F_{\mathcal{H}}(V).
\end{equation}

由于 $\mathcal{M}(\operatorname{Aut}_ZO)^{fl}_{\mathfrak{z}}$ 的作用与 Hecke 作用
交换，我们得到典范等同
$F_{\mathcal{H}}(V_1\otimes V_2)=F_{\mathcal{H}}(V_1)\otimes F_{\mathcal{H}}(V_2)$，
这意味着
\begin{equation}\tag{286}
F_{\mathcal{H}}:\ \operatorname{Rep}{}^{L}G\to
\mathcal{M}(\operatorname{Aut}_ZO)^{fl}_{\mathfrak{z}}
\end{equation}
是幺半函子。

\bdp{5.5.2} \emph{引理.} 对任意 $V\in\operatorname{Rep}{}^{L}G$，自由
$\mathfrak{z}$-模 $F_{\mathcal{H}}(V)$ 有有限秩。

\emph{证明.} 由于 $F_{\mathcal{H}}$ 是幺半函子，$F_{\mathcal{H}}(V^{*})$ 在
幺半范畴的意义下对偶于 $F_{\mathcal{H}}(V)$（见 [Del91] 的 2.1.2）。若自由
$\mathfrak{z}$-模在幺半范畴的意义下有对偶，则其秩有限。\hfill$\Box$

\SourcePage{225}

设
\begin{equation}\tag{287}
F_{\mathfrak{F}_{\mathcal{L}}}:\ \operatorname{Rep}{}^{L}G\to
\mathcal{M}(\operatorname{Aut}_ZO)^{fl}_{\mathfrak{z}}
\end{equation}
为张量函子
$F_{\mathfrak{F}_{\mathcal{L}}}(V)=V_{\mathfrak{F}_{\mathcal{L}}}$（见 5.4.10）。

现在，我们的局部主定理 5.4.11 可重新表述如下。

\bdp{5.5.3} \emph{定理.} 幺半函子 $F_{\mathcal{H}}$ 与
$F_{\mathfrak{F}_{\mathcal{L}}}$ 典范同构。

我们将说明，$F_{\mathcal{H}}$ 确实来自某个由 $\operatorname{Spec}\mathfrak{z}$
参数化的、典范定义的局部 oper 族。首先核对 $F_{\mathcal{H}}$ 确实来自
$\operatorname{Spec}\mathfrak{z}$ 上的一个 ${}^{L}G$-torsor。

\bdp{5.5.4} \emph{命题.} 幺半函子 $F_{\mathcal{H}}$ 是张量函子，即它与交换
约束相容。

证明分两步。首先把相容同构
$F_{\mathcal{H}}(V_1)\otimes F_{\mathcal{H}}(V_2)\simeq
F_{\mathcal{H}}(V_1\otimes V_2)$ 写成（扭）$\mathcal{D}$-模截面的卷积积
（见 5.5.5、5.5.6）。然后利用卷积的融合图景，说明它交换（见 ???）。
% [Tutor Note: 原书 p.225 此处交叉引用印作「???」，底本未补，译文照录。]

\bdp{5.5.5} 我们借助 Satake 等价 $h$（见（273）），把 ${}^{L}G$-模的张量范畴
换成仿射 Grassmann 流形上 $\mathcal{D}$-模的张量范畴。对
$P\in\mathcal{P}^{\natural}$，置 $F_{\mathcal{H}}(P):=F_{\mathcal{H}}(hP)$。
于是（见（282））
\begin{equation}\tag{288}
F_{\mathcal{H}}(P)=\Gamma(\mathcal{GR},P\lambda_{\mathcal{L}}^{-1})^{G(O)}.
\end{equation}

\emph{注.} 回忆 $P$ 是「超」$\mathcal{D}$-模，$\lambda_{\mathcal{L}}$ 是「超」
线丛。然而它们的奇偶性相同（都等于 $\mathcal{GR}$ 各连通分支的奇偶性），故
$P\lambda_{\mathcal{L}}^{-1}$ 是通常的偶层。这些「超」的微妙之处在我们转向
交换约束时才会起作用。

为描述相容同构
$F_{\mathcal{H}}(P_1)\otimes F_{\mathcal{H}}(P_2)\simeq
F_{\mathcal{H}}(P_1\circledast P_2)$，考虑 $\mathcal{O}^{!}$-模的积分态射
（我们用 5.3.5 的记号；关于积分见 7.11.16 的（??））
\begin{equation}\tag{289}
i_m:\ m\cdot(P_1\boxtimes'P_2)\to P_1\circledast P_2.
\end{equation}
% [Tutor Note: 原书 p.225 印作「7.11.16 (??)」，括号内编号未补全，译文照录。]

\SourcePage{226}

$\mathcal{GR}$ 上的线丛 $\lambda_{\mathcal{L}}$ 是 $G(O)$-等变的，它经由
$m:G(K)\times_{G(O)}\mathcal{GR}\to\mathcal{GR}$ 的拉回被典范地等同于「扭积」
$\lambda_{\mathcal{L}}\boxtimes'\lambda_{\mathcal{L}}$\footnote{这是因为按定义，
$\lambda_{\mathcal{L}}$ 来自 $G(K)$ 的一个中心扩张，该扩张在 $G(O)$ 上有分裂。}。
于是用 $\lambda_{\mathcal{L}}$ 扭 $i_m$，得到态射
\[
m\cdot((P_1\lambda_{\mathcal{L}}^{-1})\boxtimes'(P_2\lambda_{\mathcal{L}}^{-1}))
\to(P_1\circledast P_2)\lambda_{\mathcal{L}}^{-1}.
\]

过渡到 $G(O)$-不变截面，就得到\emph{卷积映射}（注意 $G(O)$-不变性允许我们
忽略扭）
\begin{equation}\tag{290}
*:\ \Gamma(\mathcal{GR},P_1\lambda_{\mathcal{L}}^{-1})^{G(O)}\otimes
\Gamma(\mathcal{GR},P_2\lambda_{\mathcal{L}}^{-1})^{G(O)}\to
\Gamma(\mathcal{GR},(P_1\circledast P_2)\lambda_{\mathcal{L}}^{-1})^{G(O)}.
\end{equation}

\bdp{5.5.6} \emph{引理.} 卷积映射与相容同构
$F_{\mathcal{H}}(P_1)\otimes F_{\mathcal{H}}(P_2)\simeq
F_{\mathcal{H}}(P_1\circledast P_2)$ 一致。

\emph{证明.} 考虑典范同构（作用约束）
$a:P_1\circledast(P_2\circledast\mathit{Vac}')\xrightarrow{\sim}
(P_1\circledast P_2)\circledast\mathit{Vac}'$。对
$f\in\operatorname{Hom}(\mathit{Vac}',P_1\circledast\mathit{Vac}')$、
$g\in\operatorname{Hom}(\mathit{Vac}',P_2\circledast\mathit{Vac}')$，相容同构
把 $f\otimes g$ 送到 $(P_1\circledast g)\circ f$。\hfill$\Box$

\bdsec{5.6. 重整化泛包络代数}

\bdp{5.6.1} 设 $A$ 为 $\widetilde{\mathfrak{g}\otimes K}$ 的完备化泛包络代数。
由 3.6.2，$A$ 是 $\mathbb{C}[h]$ 上的平坦代数，$h:=\mathbf{1}-1$，且
$A/hA=\overline{U}'$。

$A$ 上的自然拓扑在
$A[h^{-1}]:=A\otimes_{\mathbb{C}[h]}\mathbb{C}[h,h^{-1}]$ 上诱导一个拓扑；
事实上这是归纳极限拓扑（把 $A[h^{-1}]$ 表示成 $A\to A\to\ldots$ 的归纳极限，
其中每个箭头都是乘以 $h$）。

设 $I\subset\mathfrak{Z}$ 是 3.6.5 中的理想。记 $J$ 为
$I\overline{U}'\subset\overline{U}'=A/hA$ 在 $A$ 中的原像（$I\overline{U}'$ 按
拓扑意义理解，即 $I\overline{U}'$ 是由 $I$ 生成的 $\overline{U}'$ 的闭理想）。
$J$ 是 $A$ 中包含 $hA$ 的闭理想。记 $A^{\natural}$ 为递增序列
$A\subset h^{-1}J\subset h^{-2}J^{2}\subset\ldots$ 的并，其中 $J^{k}$ 按拓扑
意义理解。最后置 $U^{\natural}:=A^{\natural}/hA^{\natural}$。

\SourcePage{227}

$A^{\natural}$ 是 $\mathbb{C}[h]$ 上的拓扑代数（拓扑由 $A[h^{-1}]$ 诱导）。故
$U^{\natural}$ 是拓扑 $\mathbb{C}$-代数（$U^{\natural}$ 带商拓扑）。

\bdp{5.6.2} 置 $\mathit{Vac}_A=A/A(\mathfrak{g}\otimes O)$，其中
$A(\mathfrak{g}\otimes O)$ 表示 $A$ 中由 $\mathfrak{g}\otimes O$ 生成的闭左理想。
$I$ 在 $\mathit{Vac}'=\mathit{Vac}_A/h\mathit{Vac}_A$ 上平凡地作用。由于
$\mathit{Vac}_A$ 是平坦 $\mathbb{C}[h]$-模，$A^{\natural}$ 作用在
$\mathit{Vac}_A$ 上。因此 $U^{\natural}$ 作用在 $\mathit{Vac}'$ 上。

\bdp{5.6.3} 记 $U^{\natural}_0$ 为 $A$ 在 $U^{\natural}$ 中的像。
$U^{\natural}_0$ 是 $U^{\natural}$ 的子代数。我们给 $U^{\natural}_0$ 配上诱导
拓扑。映射 $A\to U^{\natural}_0$ 经由 $A/hA=\overline{U}'$ 分解，事实上经由
$\overline{U}'/I\overline{U}'$ 分解。于是得到满的连续同态
$f:\overline{U}'/I\overline{U}'\to U^{\natural}_0$。$f$ 大概是同胚。

无论如何，$f$ 诱导拓扑同构
$\mathfrak{z}=\mathfrak{Z}/I\xrightarrow{\sim}f(\mathfrak{z})$（用
$U^{\natural}$ 在 $\mathit{Vac}'$ 上的作用）。我们将把 $\mathfrak{z}$ 与
$f(\mathfrak{z})$ 等同。

\bdp{5.6.4} 设 $J_I\subset A$ 表示 $I\subset\overline{U}'=A/hA$ 的原像。记
$U^{\natural}_1$ 为 $h^{-1}J_I$ 在 $U^{\natural}$ 中的像。给 $U^{\natural}_1$ 配上
由 $U^{\natural}$ 诱导的拓扑。拓扑代数 $U^{\natural}$ 由 $U^{\natural}_1$ 生成。

\bdp{5.6.5} \emph{引理.}
\begin{enumerate}[label=(\roman*),leftmargin=2.2em]
\item $U^{\natural}_1$ 是 $U^{\natural}$ 的李子代数。
\item $U^{\natural}_0$ 是 $U^{\natural}_1$ 的理想。
\item $\mathfrak{z}U^{\natural}_1\subset U^{\natural}_1$，
$U^{\natural}_1\mathfrak{z}\subset U^{\natural}_1$。
\item $[U^{\natural}_1,\mathfrak{z}]\subset\mathfrak{z}$。
\end{enumerate}

\emph{证明.} 我们要用到 3.6.2 中定义的 Hayashi 括号 $\{\ ,\ \}$ 的若干性质。
(i) 由包含 $[J_I,J_I]\subset hJ_I$ 推出，后者显然，因为 $\{I,I\}\subset I$
（见 3.6.4 (i)）。(ii) 与 (iii) 显然。(iv) 显然，因为
$\{I,\mathfrak{Z}\}\subset\{\mathfrak{Z},\mathfrak{Z}\}\subset\mathfrak{Z}$。
\hfill$\Box$

\bdp{5.6.6} 由 5.6.5 可知，$U^{\natural}_1/U^{\natural}_0$ 是 $\mathfrak{z}$ 上的
拓扑李代数胚。乘以 $h^{-1}$ 定义了映射 $J_I\to A^{\natural}$，它诱导李代数胚
态射
\begin{equation}\tag{291}
I/I^{2}=J_I/(J_I^{2}+hA)\to U^{\natural}_1/U^{\natural}_0
\end{equation}

\SourcePage{228}

（$I/I^{2}$ 上代数胚结构的定义见 3.6.5。）态射（291）连续且满。事实上它是
拓扑同构（见 ???）。
% [Tutor Note: 原书 p.228 此处交叉引用印作「???」，底本未补，译文照录。]

\bdp{5.6.7} 记 $U^{\flat}_i$ 为 $U^{\natural}_i$ 中消灭 $\mathit{Vac}'$ 真空向量的
元素之集，$i=0,1$。引理 5.6.5 在把 $U^{\natural}_i$ 换成 $U^{\flat}_i$
（$i=0,1$）后仍然成立。故 $U^{\flat}_1/U^{\flat}_0$ 是 $\mathfrak{z}$ 上的拓扑
李代数胚。自然映射 $U^{\flat}_1/U^{\flat}_0\to U^{\natural}_1/U^{\natural}_0$ 是
拓扑同构。于是（291）诱导满的连续李代数胚态射
\begin{equation}\tag{292}
I/I^{2}\to U^{\flat}_1/U^{\flat}_0.
\end{equation}

\bdp{5.6.8} 设 $V$ 是拓扑 $U^{\natural}$-模（在我们关心的应用中 $V$ 是离散的）。
则 $V^{\mathfrak{g}\otimes O}$ 是李代数胚 $I/I^{2}$ 上的（左）拓扑模。事实上，
首先 $V^{\mathfrak{g}\otimes O}$ 是 $\mathfrak{z}$-模。其次，
$V^{\mathfrak{g}\otimes O}=\{v\in V\mid U^{\flat}_0v=0\}$，故李代数
$U^{\flat}_1/U^{\flat}_0$ 作用在 $V^{\mathfrak{g}\otimes O}$ 上。若
$v\in V^{\mathfrak{g}\otimes O}$、$z\in\mathfrak{z}$、
$a\in U^{\flat}_1/U^{\flat}_0$，则 $a(zv)-z(av)=\partial_a(z)v$，其中
$\partial_a\in\operatorname{Der}\mathfrak{z}$ 按 $U^{\flat}_1/U^{\flat}_0$ 上的
代数胚结构对应于 $a$。于是 $V^{\mathfrak{g}\otimes O}$ 是代数胚
$U^{\flat}_1/U^{\flat}_0$ 上的模。利用（292）可见 $V^{\mathfrak{g}\otimes O}$ 是
李代数胚 $I/I^{2}$ 上的模。

\bdp{5.6.9} 由（89），有连续李代数态射
$\operatorname{Der}O\to h^{-1}J_I\subset A[h^{-1}]$，使得
$L_n\mapsto h^{-1}\widetilde{\mathcal{L}}_n$，$n\ge -1$。它诱导连续李代数态射
$\operatorname{Der}O\to U^{\flat}_1\subset U^{\natural}$。另一方面，我们在
3.6.16 中定义了典范态射 $\operatorname{Der}O\to I/I^{2}$。显然下图交换：
\[
\begin{array}{ccc}
\operatorname{Der}O & \longrightarrow & U^{\flat}_1\\[2pt]
\downarrow & & \downarrow\\[2pt]
I/I^{2} & \longrightarrow & U^{\flat}_1/U^{\flat}_0
\end{array}
\]

\emph{注.} 态射 $\operatorname{Der}O\to U^{\flat}_1/U^{\flat}_0$ 把
$\operatorname{Der}O$ 同胚地映到其像上。由于 $U^{\flat}_1/U^{\flat}_0$ 连续
作用在 $\mathfrak{z}\subset U^{\natural}_0$ 上，这可由
\SourcePage{229}
关于态射 $\operatorname{Der}O\to\operatorname{Der}\mathfrak{z}$ 的相应论断推出，
而后者显然（看 $\mathfrak{z}$ 的 Sugawara 元素）。

\bdp{5.6.10} 设我们处于 5.6.8 的情形。由 5.6.9，$\operatorname{Der}O$ 经由
态射 $\operatorname{Der}O\to U^{\natural}$ 作用在 $V$ 上，子空间
$V^{\mathfrak{g}\otimes O}$ 是 $\operatorname{Der}O$-不变的，且
$\operatorname{Der}O$ 在 $V^{\mathfrak{g}\otimes O}$ 上的作用与来自态射
$\operatorname{Der}O\to I/I^{2}$ 的作用一致。

\bdp{5.6.11} \emph{注.} 2.5.1 中 $\widetilde{\mathfrak{g}\otimes K}$ 的定义涉及
由（18）定义的「临界」数量积 $c$。设想考虑对应于 $\lambda c$ 的中心扩张
$0\to\mathbb{C}\to(\widetilde{\mathfrak{g}\otimes K})_{\lambda}\to
\mathfrak{g}\otimes K\to 0$，$\lambda\in\mathbb{C}^{*}$。记 $A_{\lambda}$ 为
$(\widetilde{\mathfrak{g}\otimes K})_{\lambda}$ 的完备化泛包络代数。若把 $A$ 与
$h=\mathbf{1}-1$ 换成 $A_{\lambda}$ 与
$h_{\lambda}:=\mathbf{1}_{\lambda}-\lambda^{-1}$，则 $U^{\natural}$ 的构造与映射
（291）仍然成立；这里 $\mathbf{1}_{\lambda}$ 表示
$1\in\mathbb{C}\subset(\widetilde{\mathfrak{g}\otimes K})_{\lambda}$。记
$U^{\natural}_{\lambda}$ 与 $f_{\lambda}$ 为对应于 $\lambda$ 的
$U^{\natural}$ 与（291）的类似物。可以利用典范同构
$\widetilde{\mathfrak{g}\otimes K}\xrightarrow{\sim}
(\widetilde{\mathfrak{g}\otimes K})_{\lambda}$（它使
$\mathbf{1}\mapsto\lambda\cdot\mathbf{1}_{\lambda}$）把 $A_{\lambda}$ 与
$U^{\natural}_{\lambda}$ 分别与 $A$ 与 $U^{\natural}$ 等同。此时 $f_{\lambda}$
\emph{确实}依赖于 $\lambda$：事实上 $f_{\lambda}=\lambda f_1$。

% 尾部接 p.230（本页末句完结于 f_λ = λ f_1，5.6 节尚未结束）。
